\documentclass[10pt,twoside]{article}
\usepackage[T1]{fontenc}
\usepackage[utf8]{inputenc}
\usepackage{times,microtype}
\usepackage[papersize={16.5cm,23.5cm},top=1.5cm,bottom=1.3cm,left=1.5cm,right=1.5cm]{geometry}
\usepackage{fancyhdr}
\pagestyle{fancy}\fancyhf{}
\renewcommand{\headrulewidth}{0.2pt}
\fancyhead[LE,RO]{\thepage}
\fancyhead[RE]{\footnotesize Abstract Booklet}
\fancyhead[LO]{\footnotesize ISC18}
\begin{document}
\begin{center}
{\Large\bfseries 18th Iranian Statistical Conference}\\[0.3em]
{\large Abstract Booklet (English)}\\[0.3em]
{\small\itshape Sorted by last name of first author}
\end{center}
\newpage

% p.1243
\noindent\begin{minipage}[t][\textheight]{\textwidth}
\noindent{\bfseries\large Graph Feature Extraction for Path Planning Using Machine Learning by SVM and GNN}\\[0.35em]
{\itshape Mohsen Abdolhosseinzadeh}\\[0.2em]
{\footnotesize First author: Mohsen Abdolhosseinzadeh}\\[0.45em]
\noindent\rule{\textwidth}{0.3pt}\\[0.4em]
{\footnotesize  This paper surveys how machine learning, particularly graph neural networks and transformer architectures, automates the extraction of features to improve path planning processes. We characterize the methodological development from traditional topological invariant graph features to advanced data-driven representation learning frameworks, demonstrating their efficacy in optimizing navigation policies across complex, stochastic, and multi-agent operational domains. Furthermore, the practical realities of implementing these models are addressed, focusing on computational limits at the edge, real-world execution, and enhancing models in operational settings. Finally, experimental results are provided supporting the efficacy of the proposed graph feature extraction framework in reducing node expansions and improving path optimality across a variety of path planning benchmarks.}\\[0.4em]
{\footnotesize\bfseries Keywords:} {\footnotesize Path planning, Graph Neural Networks, Heuristic learning, Problem reduction. \textbackslash\{\}\textbackslash\{\}}\\[0.15em]
{\footnotesize\bfseries MSC:} {\footnotesize 68T20 , 68T05 , 90C27 . 1cm}\\
\vfill
\end{minipage}
\newpage

% p.1191
\noindent\begin{minipage}[t][\textheight]{\textwidth}
\noindent{\bfseries\large Analysis of social network Facebook comments using ML, Deep Learning, and lexicon-based}\\[0.35em]
{\itshape Sheida Abdolmaleki, Vadood Keramati Sh}\\[0.2em]
{\footnotesize First author: Sheida Abdolmaleki}\\[0.45em]
\noindent\rule{\textwidth}{0.3pt}\\[0.4em]
{\footnotesize  Customer satisfaction is the most important tool for understanding how our application works. Facebook is one of the most popular applications in the world that most people use. Hence, analyzing the comments in the application is useful for knowing its performance. The analysis was conducted using a series of specialized functions commonly employed in these methods to evaluate people's opinions and sentiments toward the program. The opinions and sentiments are classified into 4 classes (natural, positive, negative, and empty). In the end, the results of the two methods were compared to determine which method performed better for data analysis.}\\[0.4em]
{\footnotesize\bfseries Keywords:} {\footnotesize Sentiment Analysis, Machine learning, Deep learning, Lexicon-based \textbackslash\{\}\textbackslash\{\}}\\[0.15em]
{\footnotesize\bfseries MSC:} {\footnotesize 99X99 , 99X99 , 99X99 . 1cm}\\
\vfill
\end{minipage}
\newpage

% p.1181
\noindent\begin{minipage}[t][\textheight]{\textwidth}
\noindent{\bfseries\large Enhancing Decision Tree Learning via Conditional Cumulative Residual Entropy: A Framework for Continuous Variables}\\[0.35em]
{\itshape Somayeh Abolhosseini, Mohammad Khorashadizadeh}\\[0.2em]
{\footnotesize First author: Somayeh Abolhosseini}\\[0.45em]
\noindent\rule{\textwidth}{0.3pt}\\[0.4em]
{\footnotesize  Conventional decision tree algorithms struggle with continuous variables, often causing information loss through discretization. This paper proposes a novel framework using Conditional Cumulative Residual Entropy (CCRE) to directly handle continuous inputs and outputs. Unlike traditional entropy measures, CCRE provides a more nuanced uncertainty measure for continuous distributions. The theoretical framework for the splitting criterion is derived, and its effectiveness is shown across regression and classification tasks. Empirical evaluations indicate that this approach eliminates discretization biases and yields superior predictive performance compared to baseline methods. The results highlight CCRE's potential as a powerful tool for tree-based learning, contributing significantly to data mining and machine learning.\textbackslash\{\}\textbackslash\{\}}\\[0.4em]
{\footnotesize\bfseries Keywords:} {\footnotesize Conditional Cumulative Residual Entropy, Decision tree, Machin learning \textbackslash\{\}\textbackslash\{\}}\\[0.15em]
{\footnotesize\bfseries MSC:} {\footnotesize 99X99 , 99X99 , 99X99 . 1cm}\\
\vfill
\end{minipage}
\newpage

% p.1216
\noindent\begin{minipage}[t][\textheight]{\textwidth}
\noindent{\bfseries\large Multilevel Analysis of Misclassified Binary Outcomes: Validation-Data-Based Approach Adjusting for Covariate Measurement Error}\\[0.35em]
{\itshape Maryam Ahangari}\\[0.2em]
{\footnotesize First author: Maryam Ahangari}\\[0.45em]
\noindent\rule{\textwidth}{0.3pt}\\[0.4em]
{\footnotesize  Longitudinal designs inherently involve correlated repeated measurements nested within individuals. In such frameworks, time-varying covariates and discrete response variables are frequently subject to error due to methodological constraints or prohibitive costs. Naive estimation strategies that ignore errors in either the covariates or the response, systematically yield biased parameter estimates, affecting the validity of statistical inference and reducing asymptotic efficiency. To address this, we propose a joint modeling framework that simultaneously adjusts for misclassification in correlated binary responses and measurement error in continuous covariates. Our methodology leverages validation data to reliably estimate the misclassification mechanism and employs multivariate Gauss-Hermite quadrature to approximate the intractable integrals within the marginal likelihood function. Simulation studies are conducted to evaluate the finite-sample performance of the proposed estimator.}\\[0.4em]
{\footnotesize\bfseries Keywords:} {\footnotesize Measurement Error, Internal Validation Data, Multivariate Gauss-Hermite Quadrature, Random Effects Logistic Regression Model, Misclassification. \textbackslash\{\}\textbackslash\{\}}\\[0.15em]
{\footnotesize\bfseries MSC:} {\footnotesize 62J12 , 62F10 . 1cm}\\
\vfill
\end{minipage}
\newpage

% p.1268
\noindent\begin{minipage}[t][\textheight]{\textwidth}
\noindent{\bfseries\large Grey Relational Analysis: A Statistical Approach to Measuring Dependence in Data-Poor Conditions}\\[0.35em]
{\itshape Mohamad Ahmadian}\\[0.2em]
{\footnotesize First author: Mohamad Ahmadian}\\[0.45em]
\noindent\rule{\textwidth}{0.3pt}\\[0.4em]
{\footnotesize Grey relational analysis is one of the most important tools in grey systems theory that is used to measure the intensity of the relationship between variables in situations where data are limited, incomplete, or uncertain. In terms of functionality, this method is similar to some statistical methods such as correlation and dependence analysis, but unlike classical statistics, it does not depend on strict assumptions such as normality of distribution, large sample size, or linearity of the relationship. Grey relational analysis, by focusing on the similarity of the pattern of changes between data series, allows ranking of effective factors and identifying important variables. In this article, the theoretical foundations, formulation, implementation steps, its connection with statistical concepts, and its applications in data analysis are reviewed. Keywords: Grey relational analysis, grey statistics, statistical dependence, correlation, incomplete data, factor ranking enumerate Introduction enumerate In many statistical studies, the main goal is to examine the relationship between variables. Researchers usually want to know which factor has the greatest impact on an indicator, which variables are dependent on each other, and to what extent this dependence is. In classical statistics, tools such as Pearson's correlation coefficient, Spearman's correlation, regression, and analysis of variance are used to answer these questions. However, these methods often depend on conditions such as a sufficient number of data, appropriate quality of observations, and the establishment of some statistical assumptions (such as normality of distribution). In real problems, especially in the humanities, management, economics, environmental science, and some engineering fields, data are often limited and incomplete. In such situations, Grey Relational Analysis (GRA) is proposed as a complementary and, in some cases, alternative method for measuring the dependence between variables. This method, which belongs to the field of ``Grey Systems Theory'', is designed for systems with incomplete information or high uncertainty (Ahmadian, 2026). Unlike classical methods based on probability distributions, GRA is based on the ``geometric similarity'' between data sequences. Simply put, this method examines whether the trend of changes in one variable (reference sequence) is consistent with the trend of changes in other variables (comparison sequences). The closer the shape of the data curves, the stronger the grey relationship and the higher the grey correlation coefficient (Shi et al., 2026). The main advantage of this method is the lack of strict statistical assumptions and high efficiency with small data volumes. This feature has made GRA an ideal tool for decision-making in complex and uncertain situations. In the following, new applications of this method in scientific articles are reviewed. enumerate Basics of Grey Relational Analysis and its Relationship to Statistics enumerate This section provides a precise definition of gray relational analysis and its place in statistics. It then examines how this method fits into the framework of statistical analysis and what role it plays in complementing classical methods. 2.1 Definition and Concept of Grey Relational Analysis Grey relational analysis is a method for measuring the degree of closeness or similarity between a reference sequence and several comparison sequences (Liu et al., 2022). If the changes of a variable during observations have a pattern similar to the changes of the target variable, it is said that the degree of Grey relational of that variable with the target variable is high. As a result, that variable is considered a more important factor in explaining the behavior of the dependent variable from a statistical and analytical point of view. From a statistical perspective, Grey relational analysis is a type of non-parametric dependence measurement. In this method, instead of estimating distribution parameters or calculating covariance, the relative distance between normalized data is used. Therefore, this method is very suitable for data that: have a small sample size, do not follow a specific distribution, have different scales, have irregularity and incomplete information. 2.2 The Place of Grey Relational Analysis in Statistics Grey relational analysis is similar in purpose to statistical methods of measuring relationships, but it differs from them in basis. In statistics, correlation indicates the degree of covariation between two variables. For example, Pearson's correlation coefficient focuses on the linear relationship and numerical structure of the data. In contrast, Grey relational analysis focuses on the geometric similarity and pattern of changes in sequences. In statistical terms: If correlation is concerned with ``quantitative covariance'', Grey relational analysis is concerned with ``co-direction and co-formity of trends''. Therefore, two variables may not have a high correlation coefficient, but they are close to each other in terms of Grey relational analysis, because their pattern of increase and decrease is similar. This feature has made Grey relational analysis a useful tool for exploratory data analysis, variable selection, and factor ranking. Grey relational analysis is often compared to the correlation coefficient, but there are important differences between them. First, Pearson correlation measures the strength of a linear relationship between two variables, while the Grey relational analysis does not depend on linearity. Second, classical correlation usually requires a larger sample size for more robust results, while the Grey relational analysis can also be applied to small samples. Third, in correlation, the presence of outliers can have a large effect on the result, but in Grey relational analysis, due to the emphasis on the pattern of change, this effect is in some cases less. However, it should be noted that the Grey relational analysis is not a complete substitute for correlation. Correlation is an accurate statistical indicator for measuring dependence in a probability framework, but the Grey relational analysis is more of a tool for analyzing poorly informed systems and ranking factors in conditions of insufficient data. enumerate Grey Relational Analysis Method enumerate The calculation method for GRA is as follows (Tharisanan et al., 2026): Suppose a reference sequence is available as Eq.1. tabular |p 0.9in |p 2.6in |p 0.9in | \& X\_ 0 =\textbackslash\{\} x\_ 0 GrindEQ\_\_1\_ ,x\_ 0 GrindEQ\_\_2\_ ,...,x\_ 0 (n)\textbackslash\{\} \& tabular m comparison sequences are also defined as Eq.2. tabular |p 0.9in |p 2.6in |p 0.9in | \& X\_ i =\textbackslash\{\} x\_ i GrindEQ\_\_1\_ ,x\_ i GrindEQ\_\_2\_ ,...,x\_ i (n)\textbackslash\{\} , i=1,2,...,m \& GrindEQ\_\_2\_ tabular The goal is to calculate the degree of relational of each comparison sequence to the reference sequence. Since the variables may be measured on different scales, the data must first be normalized. This step is very important in statistics, because direct comparison of dissimilar data can be misleading. Depending on the nature of the indicator, different methods of scale-independence are used; such as: initial normalization, conversion to mean ratio, and scaling between zero and one. Normalization allows the analysis to focus on the pattern of changes, rather than the absolute magnitude of the numbers. After normalization, the absolute difference between the reference sequence and each comparison sequence is calculated for each observation is available as Eq.3. tabular |p 0.9in |p 2.6in |p 0.9in | \& \_ i (k)=x\_ 0 (k)-x\_ i (k) \& GrindEQ\_\_3\_ tabular This difference is statistically a measure of the distance between corresponding observations. Then, in the entire data, the smallest and largest difference values are determined as Eq.4 and Eq.5. tabular |p 0.9in |p 2.6in |p 0.9in | \& \_ = \_ i \_ k \_ i (k) \& GrindEQ\_\_4\_ \& \_ = \_ i \_ k \_ i (k) \& GrindEQ\_\_5\_ tabular These two values are used to scale the distances and calculate the relational coefficient. The Grey relational coefficient for sequence i at point k is obtained from Eq.6. tabular |p 0.9in |p 2.6in |p 0.9in | \& \_ i (k)= \_ + \_ \_ i (k)+ \_ \& GrindEQ\_\_6\_ tabular In this relationship, is the discrimination coefficient and its value is usually taken to be 0.5. This parameter plays the role of regulating the sensitivity of the model from a statistical point of view. The smaller is, the greater the discrimination between the sequences. The value of i(k) is between zero and one. The closer this value is to one, the greater the similarity of that observation from the comparison sequence to the reference sequence. To obtain an overall index of the degree of dependence of each variable with the reference variable, the average of the relational coefficients is calculated in Eq.7. tabular |p 0.9in |p 2.6in |p 0.9in | \& \_ i = 1 n \_ k=1 \textasciicircum{} n \_ i (k) \& GrindEQ\_\_7\_ tabular The Grey relational degree i is actually a summary statistic that shows the intensity of the overall relationship between the two sequences. The larger this value, the stronger the statistical relationship between the variable of interest and the reference variable is considered. enumerate Example of Grey Relational Analysis enumerate To illustrate the application of grey relational analysis, consider a bank that seeks to identify which operational indicator is most closely associated with the monthly profit of a branch. Three explanatory variables are considered: the number of active customers (X \_ 1 ), the number of digital transactions (X \_ 2 ), and the number of customer complaints (X \_ 3 ). The response variable (X \_ 0 ) is the monthly profit (thousand USD). Table 1 presents a simple illustrative dataset. Table 1 Sample data tabular |p 0.7in |p 1.0in |p 0.9in |p 0.9in |p 0.9in | Month \& X \_ 0 \& X \_ 1 \& X \_ 2 \& X \_ 3 1 \& 100 \& 96 \& 99 \& 80 2 \& 120 \& 115 \& 120 \& 95 3 \& 140 \& 133 \& 141 \& 120 tabular The absolute differences are shown in Table 2. Table 2 Absolute difference sequences tabular |p 0.7in |p 1.3in |p 1.2in |p 1.2in | Month \& 1 \& 2 \& 3 1 \& 4 \& 1 \& 20 2 \& 5 \& 0 \& 25 3 \& 7 \& 1 \& 20 tabular Thus, \_ =0 , \_ =25 The grey relational coefficients, calculated using the commonly adopted distinguishing coefficient =0.5, are presented in Table 3. Table 3 Grey relational coefficients tabular |p 0.7in |p 1.3in |p 1.2in |p 1.2in | Month \& 1 \& 2 \& 3 1 \& 0.758 \& 0.926 \& 0.385 2 \& 0.714 \& 1.000 \& 0.333 3 \& 0.641 \& 0.926 \& 0.385 tabular The grey relational grades are obtained by averaging the coefficients. Accordingly, 1 =0.704, 2=0.951, 3=0.368, hence, 2 > 1 > 3. The results indicate that X \_ 2\textbackslash\{\} (digital transactions) exhibit the strongest grey relational grade with monthly profit, followed by the number of active customers, whereas customer complaints have the weakest relationship. This outcome is consistent with managerial expectations, as higher digital transaction volumes are generally associated with increased banking activity and revenue generation, while customer complaints are less directly related to profitability. Consequently, grey relational analysis provides a practical tool for ranking explanatory variables according to their association with a response variable, thereby supporting feature selection, dimensionality reduction, and sensitivity analysis. enumerate Applications, Advantages, and Limitations enumerate In this section, the practical applications of gray relational analysis in various domains are reviewed. The advantages and limitations of this method are also analyzed in detail to provide a complete view of its capabilities and challenges. 5.1 Applications This method is used in various fields for data analysis. In economics, it can be used to determine the factors affecting economic growth or inflation. In management, it is used to rank organizational performance indicators. In social sciences, it is used to measure the proximity of variables such as satisfaction, development, or quality of life. In engineering and environmental science, it is also used to evaluate parameters affecting quality, sustainability, or pollution. In all these cases, Grey relational analysis plays a role similar to a statistical tool to discover the structure of dependence between variables. For example, articles can be reviewed in the following areas: Economic studies that have used Grey relational analysis to identify factors affecting economic growth, inflation, productivity, or financial markets. Management research that has used this method to evaluate organizational performance, rank productivity indicators, select suppliers, or analyze customer satisfaction. Environmental articles that have analyzed the relationship between pollution indicators, energy consumption, water quality, or sustainable development using this method. Engineering and industrial research that has used Grey relational analysis to select optimal parameters, quality control, or assess reliability. Health and social science studies in which this method has been used to rank risk factors, analyze well-being indicators, or examine variables affecting quality of life. 5.2 Advantages of Grey relational analysis in statistical research Grey relational analysis has several important advantages that make it attractive for applied statistical studies: Requirement of small sample size, no strong dependence on distributional assumptions, ability to use incomplete or irregular data, possibility of comparing non-scale variables after normalization, suitable for ranking effective factors, relative simplicity in calculation and interpretation. For this reason, this method is widely used in research where data are limited or conditions for performing classical tests are not available. 5.3 Limitations of the Method Despite its many advantages, Grey relational analysis also has some limitations. Its results may be sensitive to the normalization method. Also, the choice of the discrimination coefficient can affect the final ranking. On the other hand, this method mostly shows the intensity of similarity and, like regression models, does not provide an explicit causal or functional relationship. As a result, from a statistical perspective, Grey relational analysis is better used as a complementary method alongside other methods such as correlation, regression, factor analysis, or multi-criteria decision-making methods. enumerate Review of recent applied studies enumerate In this section, recent research that has used GRA in diverse areas such as managing large projects, selecting sustainable suppliers in the electric vehicle supply chain, and accurately predicting fuel properties is reviewed. These studies show how combining GRA with other methods (such as fuzzy evaluation or machine learning) can be a powerful tool for decision-making under conditions of uncertainty and high complexity. Anjum et al. (2026) used Grey relational analysis to rank AI-based mental health interventions. By collecting opinions from experts and patients, it was shown that ``therapeutic effectiveness'' was the most important criterion for selecting these tools, and this priority did not differ between different demographic groups. It was also shown that the grey model can prioritize complex criteria well when compared to other methods. Xiao et al. (2026) used a dynamic lag Grey relational analysis model to show that ``new productivity'' (such as Internet development and technological innovation) reduces rural energy poverty with a time lag. It was proved that if time and lag effects are considered, the grey model can show the relationship between technology and energy poverty reduction much more accurately than conventional methods, and its effect is greater in less developed areas. Jin (2026) combined fuzzy evaluation and Grey relational analysis to provide a model to quantitatively measure the complexity of large construction projects. Using the Grey method, the weight and importance of each complexity factor (such as technical, organizational, and environmental) were calculated, and it was shown that this method can provide the most accurate picture of the complexity status of large projects without subjective bias. Singh \textbackslash\{\}\& Pandey (2026) used Grey relational analysis to select the best supplier in the electric vehicle supply chain with a focus on environmental sustainability. They examined 17 environmental criteria and showed that the method could rank suppliers well, and the fourth supplier was identified as the best option due to its lower resource consumption and use of new technologies. Lawal et al. (2025) used Grey relational analysis to build a hybrid machine learning model to improve the prediction of the calorific value of coal. It was shown that the individual models performed poorly when exposed to new data, but by combining them using the grey method, the prediction accuracy was significantly improved. enumerate Conclusion enumerate Grey relational analysis is an efficient method for measuring the dependence and ranking of variables in situations where data are limited, incomplete, or lack the assumptions necessary for classical statistical analyses. By focusing on the similarity of the pattern of changes between sequences, this method provides a non-parametric statistical measure of the relationship between variables. Although it cannot be a complete replacement for classical methods such as correlation and regression, it is very useful in exploratory analyses, variable selection, factor evaluation, and decision-making. Hence, Grey relational analysis can be considered a bridge between applied statistics and grey systems theory. References Ahmadian, M. (2026). Forecasting inbound tourist arrivals to Iran post-COVID-19 pandemic: a small data approach using grey models. Journal of Tourism Futures, 1--15. doi:https://doi.org/10.1108/JTF-10-2024-0219Anjum, R., Daud, S., Bhatti, G., \textbackslash\{\}\& Saddique, F. (2026). Evaluation of AI-based mental health interventions using the grey relational analysis. Grey Systems: Theory and Application, 16 GrindEQ\_\_1\_ , 75--90. doi:https://doi.org/10.1108/GS-03-2025-0031Jin, S. (2026). Measuring complexity in mega construction projects: fuzzy comprehensive evaluation and grey relational analysis. Engineering, Construction and Architectural Management, 33 GrindEQ\_\_1\_ , 284--317. doi:https://doi.org/10.1108/ECAM-07-2024-0951Lawal, A., Ajeboriogbon, A., Onifade, M., Oluwaseyi Bada, S., \textbackslash\{\}\& Mulenga, F. (2026). A novel grey relational analysis-based committee of machine learning methods for enhanced prediction of coal calorific value. Fuel, 406 Part C, 137070. doi:https://doi.org/10.1016/j.fuel.2025.137070Liu, S., Yang, Y., \textbackslash\{\}\& Forrest, J.-L. (2022). Grey Relational Analysis Models. In Grey Systems Analysis (pp. 77--124). Singapore: Springer. doi:https://doi.org/10.1007/978-981-19-6160-1\textbackslash\{\}\_5Shi, L., Luo, D., Zhang, M., \textbackslash\{\}\& Zhang, D. (2026). Research on the correlation degree between the manufacturing industry and the digital industry in the Yellow River Basin based on the grey matrix correlation model Available to Purchase. Grey Systems: Theory and Application, 1-23. doi:https://doi.org/10.1108/GS-12-2025-0194Singh, G., \textbackslash\{\}\& Pandey, A. (2026). Environmental sustainability integrated supplier selection in electric vehicle supply chains: a grey relational analysis approach. Environment, Development and Sustainability, 28, 7861--7889. doi:https://doi.org/10.1007/s10668-024-05294-xTharisanan, R., Kathiresan, M., Kathiravan, K., \textbackslash\{\}\& Priyavarshini, M. (2026). Multiobjective optimization of machining parameters on electrical discharge machine by Grey Relational Analysis coupled with Firefly Algorithm. Journal of the Chinese Institute of Engineers, 49 GrindEQ\_\_1\_ . doi:https://doi.org/10.1080/02533839.2025.2544774Xiao, Q., Gao, M., \textbackslash\{\}\& Wu, R. (2026). How can new quality productivity alleviate rural energy poverty? Evidence from dynamic lag grey relational analysis. Energy Sources, Part B: Economics, Planning, and Policy, 21 GrindEQ\_\_1\_ . doi:https://doi.org/10.1080/15567249.2025.2599195 document}\\[0.4em]
\vfill
\end{minipage}
\newpage

% p.1285
\noindent\begin{minipage}[t][\textheight]{\textwidth}
\noindent{\bfseries\large A Density-Augmented Representation Framework for Improving Clustering of Nonlinear Synthetic Data}\\[0.35em]
{\itshape Vahideh Ahrari}\\[0.2em]
{\footnotesize First author: Vahideh Ahrari}\\[0.45em]
\noindent\rule{\textwidth}{0.3pt}\\[0.4em]
{\footnotesize  Clustering performance may deteriorate when nonlinear geometry, heterogeneous local density, or outliers make the original feature space insufficiently informative. This paper proposes an Adaptive Hybrid Density-Augmented Clustering (AHDC) framework that enriches standardized original features with local density-based descriptors and controls the contribution of the density-feature block through an observation-specific adaptive weight. Experiments on five synthetic datasets over 20 runs show that AHDC achieves its largest observed gain for nonlinear structures, while its mean ARI remains close to that of raw-feature K-means in the Easy, Hard, and Outlier scenarios. Its performance decreases under strong cluster-size imbalance, where K-means Raw and Spectral clustering achieve higher mean ARI values. These findings indicate that adaptive density augmentation is beneficial under specific geometric conditions rather than uniformly superior across all clustering regimes. \textbackslash\{\}\textbackslash\{\}}\\[0.4em]
{\footnotesize\bfseries Keywords:} {\footnotesize adaptive clustering, density-aware representation, adaptive feature fusion, neighborhood-based features, nonlinear clustering. \textbackslash\{\}\textbackslash\{\}}\\[0.15em]
{\footnotesize\bfseries MSC:} {\footnotesize 99X99 , 99X99 , 99X99 . \textbackslash\{\}\textbackslash\{\}}\\
\vfill
\end{minipage}
\newpage

% p.1014
\noindent\begin{minipage}[t][\textheight]{\textwidth}
\noindent{\bfseries\large A Robust FSVM-GARCH Framework for Imbalanced Time Series Data}\\[0.35em]
{\itshape Najmeh Rashidi Alavijeh, Saeid Rezakhah}\\[0.2em]
{\footnotesize First author: Najmeh Rashidi Alavijeh}\\[0.45em]
\noindent\rule{\textwidth}{0.3pt}\\[0.4em]
{\footnotesize  This paper presents an enhanced fuzzy support vector machine (FSVM) framework tailored for financial time series classification. We address the persistent challenges of noise and class imbalance by introducing a novel integration of domain-based volatility information and adaptive fuzzy weighting. The key innovation lies in incorporating the Parkinson volatility estimator---a highly efficient and model-free variance measure---as the sole input feature, ensuring robustness in environments where high-frequency data are unavailable. Furthermore, we develop a weighted fuzzy membership mechanism (WFSVM) that assigns membership values based on inverse class frequency weighting and volatility-based distance metrics. This dual mechanism significantly enhances the classifier's performance, allowing it to outperform standard SVM and state-of-the-art imbalance-corrected models in volatility regime recognition, especially under skewed data distributions.}\\[0.4em]
{\footnotesize\bfseries Keywords:} {\footnotesize Fuzzy Support Vector Machine (FSVM), Class Imbalance Learning (CIL), Parkinson Volatility Estimator, Volatility Clustering, Mahalanobis Distance\textbackslash\{\}\textbackslash\{\}}\\[0.15em]
{\footnotesize\bfseries MSC:} {\footnotesize 62H30 , 91G70 , 68T10 . 1cm}\\
\vfill
\end{minipage}
\newpage

% p.1187
\noindent\begin{minipage}[t][\textheight]{\textwidth}
\noindent{\bfseries\large Comparative Evaluation of Regression Kriging and Smoothing Splines for Spatial Prediction of Soil Zinc Contamination}\\[0.35em]
{\itshape Roshanak Alimohammadi, Fatemeh Mahdioun}\\[0.2em]
{\footnotesize First author: Roshanak Alimohammadi}\\[0.45em]
\noindent\rule{\textwidth}{0.3pt}\\[0.4em]
{\footnotesize Accurate mapping of soil contamination is essential for environmental assessment and remediation planning. This study compares Regression Kriging (RK) and Smoothing Splines (SS) for estimating zinc (Zn) concentrations in the Meuse River Basin using distance from river as an environmental covariate. Variogram analysis indicated a strong spatial structure, with a nugget effect of 0.0356, a total sill of 0.2866, and a range of 267.3\textbackslash\{\},m. Model performance was evaluated through leave-one-out cross-validation (LOOCV). The results showed that RK outperformed SS, yielding lower RMSE (RK:0.3868 vs.\textbackslash\{\} SS:0.3989) and MAE (RK:0.2811 vs.\textbackslash\{\} SS:0.2966). By integrating regression with geostatistical modeling of residuals, RK provided more accurate predictions and better identified local contamination hotspots. These findings highlight the advantage of hybrid geostatistical approaches for spatial modeling of heavy-metal pollution in river-basin environments. 0.3cm}\\[0.4em]
{\footnotesize\bfseries Keywords:} {\footnotesize Regression Kriging, Smoothing Spline, Spatial Autocorrelation, Auxiliary Variable, Digital Soil Mapping, Soil Contamination. \textbackslash\{\}\textbackslash\{\}}\\[0.15em]
{\footnotesize\bfseries MSC:} {\footnotesize 62H11 . 1cm}\\
\vfill
\end{minipage}
\newpage

% p.1211
\noindent\begin{minipage}[t][\textheight]{\textwidth}
\noindent{\bfseries\large Periodic Splines, SARIMA, and Frequency-Domain Analysis: A Comparative Study Using the Arosa Ozone Dataset}\\[0.35em]
{\itshape Dr. Roshanak Alimohammadi, Azam Imani}\\[0.2em]
{\footnotesize First author: Roshanak Alimohammadi}\\[0.45em]
\noindent\rule{\textwidth}{0.3pt}\\[0.4em]
{\footnotesize  Modeling seasonal patterns in long-term environmental records is one of the most demanding problems. This study compares frequency-domain harmonic regression, cyclic cubic splines, and SARIMA using the Arosa monthly total column ozone series (1926--1971, n=518). Harmonic regression outperformed both alternatives across all criteria: lowest in-sample MSE (276.92), strongest cross-validation (RMSE=2.81, R \textasciicircum{}2 =0.82, MAPE=0.87\textbackslash\{\}}\\[0.4em]
{\footnotesize\bfseries Keywords:} {\footnotesize Periodic Splines, Frequency-Domain Analysis, SARIMA, The Arosa Ozone Dataset}\\[0.15em]
{\footnotesize\bfseries MSC:} {\footnotesize 62M10, 62M15, 62P12.}\\
\vfill
\end{minipage}
\newpage

% p.1086
\noindent\begin{minipage}[t][\textheight]{\textwidth}
\noindent{\bfseries\large Liu-type estimators for the Elliptical SUR Model}\\[0.35em]
{\itshape Morteza Amini, Mohammad Arashi, and Mahdi Roozbeh}\\[0.2em]
{\footnotesize First author: Morteza Amini}\\[0.45em]
\noindent\rule{\textwidth}{0.3pt}\\[0.4em]
{\footnotesize This paper investigates the application of Liu estimators within the seemingly unrelated regression (SUR) model under conditions of multicollinearity and elliptically distributed errors---an issue often observed in applied economic research. The conventional SUR setup is generalized by incorporating partially linear components, which improves adaptability for intricate econometric specifications. Asymptotic characteristics of the introduced estimators are established, alongside detailed risk assessments that highlight their theoretical benefits. A wide range of Monte Carlo experiments confirms the enhanced effectiveness of this estimator relative to traditional approaches. quotation Keywords : Elliptical errors; Multicollinearity; Partially linear model; Seemingly unrelated regression.}\\[0.4em]
\vfill
\end{minipage}
\newpage

% p.1255
\noindent\begin{minipage}[t][\textheight]{\textwidth}
\noindent{\bfseries\large Predicting Competitive Bidding Behavior Under Uncertainty Using Scenario-Informed Machine Learning}\\[0.35em]
{\itshape Mahshad Arabi, 1 , Mohammad Rahim Eivazi}\\[0.2em]
{\footnotesize First author: Mahshad Arabi}\\[0.45em]
\noindent\rule{\textwidth}{0.3pt}\\[0.4em]
{\footnotesize  Competitive bidding behavior in project tendering is difficult to predict under market volatility and uncertainty. This study proposes a scenario-informed machine learning approach to classify bidding behavior into aggressive, moderate, and conservative strategies. Using 850 completed tenders from 2018 to 2023, three algorithms, Random Forest, Support Vector Machine, and XGBoost, were trained and evaluated through a temporal train-test split. Results show that Random Forest performs best, achieving 89.2\textbackslash\{\}}\\[0.4em]
{\footnotesize\bfseries Keywords:} {\footnotesize Competitive bidding behavior; Scenario-informed machine learning; Project tendering; Uncertainty; Decision support; Random Forest. \textbackslash\{\}\textbackslash\{\}}\\[0.15em]
{\footnotesize\bfseries MSC:} {\footnotesize 62P30 , 68T01 , 90B50 . 1cm}\\
\vfill
\end{minipage}
\newpage

% p.1338
\noindent\begin{minipage}[t][\textheight]{\textwidth}
\noindent{\bfseries\large Joint Parameter Estimation in the Pareto Model under Type-II Censoring: A Bayesian Approach}\\[0.35em]
{\itshape Akbar Asgharzadeh}\\[0.2em]
{\footnotesize First author: Akbar Asgharzadeh}\\[0.45em]
\noindent\rule{\textwidth}{0.3pt}\\[0.4em]
{\footnotesize  This paper presents a Bayesian approach for joint estimation of the two-parameter Pareto model under type-II censoring. Based on censored lifetime data, the posterior distributions of the model parameters are derived by combining the likelihood function with suitable prior distributions. A two-step procedure is developed to obtain closed-form Bayesian credible regions for the parameters.}\\[0.4em]
{\footnotesize\bfseries Keywords:} {\footnotesize Pareto distribution, Bayesian credible regions, type-II censored data. \textbackslash\{\}\textbackslash\{\}}\\[0.15em]
{\footnotesize\bfseries MSC:} {\footnotesize 99X99 , 99X99 , 99X99 . 1cm}\\
\vfill
\end{minipage}
\newpage

% p.1318
\noindent\begin{minipage}[t][\textheight]{\textwidth}
\noindent{\bfseries\large Cardiovascular Diseases classification using machine learning models}\\[0.35em]
{\itshape Babak Azarnavid}\\[0.2em]
{\footnotesize First author: Babak Azarnavid}\\[0.45em]
\noindent\rule{\textwidth}{0.3pt}\\[0.4em]
{\footnotesize  Early detection of cardiovascular diseases (CVD) is vital for mitigating mortality and enhancing patient outcomes. Therefore, there is a pressing need for accurate and reliable diagnostic tools. To address this, we evaluated five machine learning models for heart disease prediction: LightGBM, XGBoost, random forest, logistic regression, and a deep neural network. We applied these models to a clinical dataset comprising 1000 patient records. Our experimental results demonstrate that LightGBM achieved the best performance, with a test accuracy of 0.995, a perfect recall score, and an F1 score of 0.996. Finally, SHAP analysis identified slope, resting blood pressure, and serum cholesterol as the most influential predictors in the dataset.}\\[0.4em]
{\footnotesize\bfseries Keywords:} {\footnotesize Cardiovascular diseases, Machine learning, Classification, LightGBM\textbackslash\{\}\textbackslash\{\}}\\[0.15em]
{\footnotesize\bfseries MSC:} {\footnotesize 68T05 , 62H30 , 62P10 . 1cm}\\
\vfill
\end{minipage}
\newpage

% p.1033
\noindent\begin{minipage}[t][\textheight]{\textwidth}
\noindent{\bfseries\large Enhanced Bivariate FGM Copulas for Accurate Dependency Analysis}\\[0.35em]
{\itshape Hakim Bekrizadeh}\\[0.2em]
{\footnotesize First author: Hakim Bekrizadeh}\\[0.45em]
\noindent\rule{\textwidth}{0.3pt}\\[0.4em]
{\footnotesize  A new class of bivariate FGM copulas is proposed, extending existing models. Its properties and association measures are derived, and applications to medical data show improved correlation and tail dependence. A simulation method is developed and validated to reproduce the original dependency structure.}\\[0.4em]
{\footnotesize\bfseries Keywords:} {\footnotesize FGM Copula, Measures of Dependence. \textbackslash\{\}\textbackslash\{\}}\\[0.15em]
{\footnotesize\bfseries MSC:} {\footnotesize 62H05 , 62H20 . 1cm}\\
\vfill
\end{minipage}
\newpage

% p.1123
\noindent\begin{minipage}[t][\textheight]{\textwidth}
\noindent{\bfseries\large Fuzzy Bayesian Inference with Soft Sets}\\[0.35em]
{\itshape M. A. Dehghanizadeh}\\[0.2em]
{\footnotesize First author: M. A. Dehghanizadeh}\\[0.45em]
\noindent\rule{\textwidth}{0.3pt}\\[0.4em]
{\footnotesize  This paper presents a novel framework for Bayesian inference designed to handle data characterized by inherent fuzziness and uncertainty, utilizing the expressive power of soft sets. We introduce the concept of a Soft Bayesian model, which extends classical Bayesian methodology by formally incorporating soft evidence and soft fuzzy prior information. The core of our approach lies in defining soft likelihoods and soft priors that can capture nuanced epistemic uncertainty. We provide the mathematical formulation for soft Bayesian updating, enabling the computation of soft posterior distributions.}\\[0.4em]
{\footnotesize\bfseries Keywords:} {\footnotesize Soft Sets, Fuzzy Sets, Bayesian Inference, Uncertainty Quantification, Soft Likelihood, Soft Prior, Fuzzy Regression, Statistical Modeling. \textbackslash\{\}\textbackslash\{\}}\\[0.15em]
{\footnotesize\bfseries MSC:} {\footnotesize 62F15 , 62F35 , 62F99 , 03E72 1cm}\\
\vfill
\end{minipage}
\newpage

% p.1121
\noindent\begin{minipage}[t][\textheight]{\textwidth}
\noindent{\bfseries\large Hypergroups-Based Modifications of a Traditional Clustering Algorithm for Heterogeneous Big Data Analytics}\\[0.35em]
{\itshape M. A. Dehghanizadeh}\\[0.2em]
{\footnotesize First author: M. A. Dehghanizadeh}\\[0.45em]
\noindent\rule{\textwidth}{0.3pt}\\[0.4em]
{\footnotesize  Classical clustering models often falter when faced with the inherent heterogeneity and uncertainty prevalent in big data. This paper introduces a novel hyperalgebraic framework, specifically leveraging hypergroups, to address these challenges by fundamentally modifying a traditional clustering algorithm. We establish a robust pipeline for representing complex, multi-valued observations within this hypergroup structure. This includes developing hyper-mean and hyper-variance estimators tailored for this new representation and proving their theoretical stability. Our framework incorporates a computationally efficient hyper-fusion algorithm designed for large-scale data, demonstrating particular effectiveness in scenarios characterized by ambiguity and non-deterministic relationships. Theoretical guarantees and simulation results confirm the superior performance of our modified hypergroup-based clustering approach over traditional methods, especially in uncertain and high-dimensional data landscapes.}\\[0.4em]
{\footnotesize\bfseries Keywords:} {\footnotesize Big data analytics, Hyperstructures, Hypergroups, Set-valued statistics, Multi-valued data, Statistical inference, Uncertainty quantification, Clustering, Simulation. \textbackslash\{\}\textbackslash\{\}}\\[0.15em]
{\footnotesize\bfseries MSC:} {\footnotesize 06F05 , 60G07 , 62H30 , 28A80 1cm}\\
\vfill
\end{minipage}
\newpage

% p.1120
\noindent\begin{minipage}[t][\textheight]{\textwidth}
\noindent{\bfseries\large Stochastic Hypergroups and Applications in Statistics}\\[0.35em]
{\itshape M. A. Dehghanizadeh}\\[0.2em]
{\footnotesize First author: M. A. Dehghanizadeh}\\[0.45em]
\noindent\rule{\textwidth}{0.3pt}\\[0.4em]
{\footnotesize  This note gives a concise but rigorous introduction to hypergroups and stochastic hypergroups, with emphasis on their relevance to statistics. We begin with the precise axiomatic definition of a hypergroup, then introduce stochastic hypergroups as convolution structures in which each convolution product is a probability measure. We present one canonical example and one highly illustrative stochastic example, and then prove three fundamental theorems concerning convolution operators, random-walk transition kernels, and statistical estimation on hypergroup-valued data. Throughout, we follow the standard framework developed in harmonic analysis and probabilistic hypergroup theory; see, for example, Jewett:1975,BloomHeyer:1995,Voit:2014 .}\\[0.4em]
{\footnotesize\bfseries Keywords:} {\footnotesize Stochastic Hypergroups, Statistical Inference on Manifolds, Harmonic Analysis, Random Walks on Algebraic Structures. \textbackslash\{\}\textbackslash\{\}}\\[0.15em]
{\footnotesize\bfseries MSC:} {\footnotesize 43A62 , 60B15 , 62H99 , 60J10 1cm}\\
\vfill
\end{minipage}
\newpage

% p.1143
\noindent\begin{minipage}[t][\textheight]{\textwidth}
\noindent{\bfseries\large A Functional Beta Regression Approach for Analyzing Bounded Data with Derivative Information}\\[0.35em]
{\itshape Hamideh Dehnavi, Davood Shahsavani}\\[0.2em]
{\footnotesize First author: Hamideh Dehnavi}\\[0.45em]
\noindent\rule{\textwidth}{0.3pt}\\[0.4em]
{\footnotesize This paper proposes a novel partial functional beta regression model for bounded responses in (0,1) . The model incorporates a functional predictor, its derivative, and additional scalar covariates, using a logit link function to model the mean of the beta distribution. We employ basis expansions and penalized maximum likelihood for estimation, capturing both global effects and dynamic behavior via the derivative term. Under standard regularity conditions, we establish the identifiability and consistency of the estimators. Simulation results and an application to near-infrared milk spectra demonstrate that integrating derivative information significantly improves predictive accuracy and provides better interpretability compared to standard functional models}\\[0.4em]
{\footnotesize\bfseries Keywords:} {\footnotesize beta regression; functional data analysis; derivative effects; penalized likelihood; NIR spectroscopy}\\[0.15em]
\vfill
\end{minipage}
\newpage

% p.1324
\noindent\begin{minipage}[t][\textheight]{\textwidth}
\noindent{\bfseries\large Two-Parameter Ridge Estimation in Seemingly Unrelated Regression Models}\\[0.35em]
{\itshape Robab Mehdizadeh Esfanjani}\\[0.2em]
{\footnotesize First author: Robab Mehdizadeh Esfanjani}\\[0.45em]
\noindent\rule{\textwidth}{0.3pt}\\[0.4em]
{\footnotesize  Seemingly Unrelated Regression (SUR) models are used when several linear regression equations are investigated simultaneously. To reduce the impact of multicollinearity in SUR models, the one-parameter ridge (Ridge-1) estimator has been proposed. As a generalization, this paper introduces the two-parameter ridge (Ridge-2) estimator for SUR models with multicollinearity. The performance of the proposed estimator is compared with Generalized Least Squares (GLS) and Ridge-1 via simulations. A real dataset on chronic renal failure is used for illustration. The results show that the Ridge-2 estimator outperforms both GLS and Ridge-1 under severe multicollinearity.}\\[0.4em]
{\footnotesize\bfseries Keywords:} {\footnotesize Multicollinearity, One-parameter ridge, Seemingly unrelated regression, Two-parameter ridge, Mean squared error. \textbackslash\{\}\textbackslash\{\}}\\[0.15em]
{\footnotesize\bfseries MSC:} {\footnotesize 62J07 , 62J05 . 1cm}\\
\vfill
\end{minipage}
\newpage

% p.1081
\noindent\begin{minipage}[t][\textheight]{\textwidth}
\noindent{\bfseries\large Kernel Density Weighted Twin SVM: A Robust Approach to Noisy Data}\\[0.35em]
{\itshape Ehsan Eshaghi, Mohammad Arashi, Jalal A. Nasiri}\\[0.2em]
{\footnotesize First author: Ehsan Eshaghi}\\[0.45em]
\noindent\rule{\textwidth}{0.3pt}\\[0.4em]
{\footnotesize  We address the sensitivity of Twin Support Vector Machine variants to label noise by proposing a density-weighted Least Squares Twin Support Vector Machine (KDEWLSTSVM). The method integrates kernel density estimation (KDE) into the least-squares optimization to assign global density-based sample weights, which naturally reduce the influence of low-density and potentially mislabeled samples while avoiding graph construction used in locally weighted approaches. Experiments on three binary classification datasets with simultaneous feature perturbations and 20\textbackslash\{\}}\\[0.4em]
{\footnotesize\bfseries Keywords:} {\footnotesize Twin support vector machine, Kernel density estimation, Label noise robustness, and Density-based weighting. \textbackslash\{\}\textbackslash\{\}}\\[0.15em]
{\footnotesize\bfseries MSC:} {\footnotesize 68T10 , 62H30 , 62G07 . 1cm}\\
\vfill
\end{minipage}
\newpage

% p.1109
\noindent\begin{minipage}[t][\textheight]{\textwidth}
\noindent{\bfseries\large On Stress-Strength Reliability for a Discrete Pareto-like Distribution}\\[0.35em]
{\itshape Davood Farbod}\\[0.2em]
{\footnotesize First author: Davood Farbod}\\[0.45em]
\noindent\rule{\textwidth}{0.3pt}\\[0.4em]
{\footnotesize  In this paper, we consider a two-parameter discrete Pareto-like distribution introduced by Kuznetsov:2001, Kuznetsov:2017 . Basic properties of the model are reviewed, and its asymptotic tail behavior is examined, confirming the Pareto-like heavy-tailed structure. The stress-strength reliability parameter R=P(X>Y) is derived when both stress and strength follow this discrete distribution. An explicit expression for the reliability measure is obtained in the general case, and a closed-form representation is derived for the identically distributed case. Numerical illustrations are presented to demonstrate the effect of the model parameters on the reliability measure.}\\[0.4em]
{\footnotesize\bfseries Keywords:} {\footnotesize discrete Pareto-like distribution, heavy-tailed distribution, Hurwitz zeta function, stress-strength reliability.\textbackslash\{\}\textbackslash\{\}}\\[0.15em]
{\footnotesize\bfseries MSC:} {\footnotesize 62E10 , 60G70 , 60E05 . 1cm}\\
\vfill
\end{minipage}
\newpage

% p.1023
\noindent\begin{minipage}[t][\textheight]{\textwidth}
\noindent{\bfseries\large On the Probability Generating Function of Discrete Distribution Generated by Cauchy Stable Density}\\[0.35em]
{\itshape Davood Farbod}\\[0.2em]
{\footnotesize First author: Davood Farbod}\\[0.45em]
\noindent\rule{\textwidth}{0.3pt}\\[0.4em]
{\footnotesize  In this paper, we study the probability generating function (PGF) of a discrete distribution induced by the Cauchy stable density, referred to as the Discrete Cauchy Distribution. An approximate analytical representation of the PGF is derived, and its asymptotic behavior is examined. The analysis shows that the PGF develops a logarithmic-type singularity in its domain, indicating that the distribution is heavy-tailed with infinite mean. These results provide a precise analytical characterization of the Discrete Cauchy Distribution and offer useful insights for modeling and inference in settings involving heavy-tailed count data.}\\[0.4em]
{\footnotesize\bfseries Keywords:} {\footnotesize Discrete Cauchy distribution, Probability generating function, Asymptotic analysis. \textbackslash\{\}\textbackslash\{\}}\\[0.15em]
{\footnotesize\bfseries MSC:} {\footnotesize 60E05 , 60E10 , 60E07 , 62E10 . 1cm}\\
\vfill
\end{minipage}
\newpage

% p.1003
\noindent\begin{minipage}[t][\textheight]{\textwidth}
\noindent{\bfseries\large Entropy in Inflated Count Data: A Comparative Analysis}\\[0.35em]
{\itshape Reza Farhadian}\\[0.2em]
{\footnotesize First author: Reza Farhadian}\\[0.45em]
\noindent\rule{\textwidth}{0.3pt}\\[0.4em]
{\footnotesize  Entropy is a fundamental measure of uncertainty in random variables and depends on the underlying probability distribution. Consequently, even small changes in the probability mass assigned to observed values may substantially affect the entropy. In this paper, we study the entropy of count data under an inflated distribution and present a theorem comparing the entropy of the inflated model with that of the corresponding regular distribution.}\\[0.4em]
{\footnotesize\bfseries Keywords:} {\footnotesize Count data, Inflated distribution, entropy. \textbackslash\{\}\textbackslash\{\}}\\[0.15em]
{\footnotesize\bfseries MSC:} {\footnotesize 94A17 , 60E05 . 1cm}\\
\vfill
\end{minipage}
\newpage

% p.1317
\noindent\begin{minipage}[t][\textheight]{\textwidth}
\noindent{\bfseries\large Analysis of Persian Legal Document Networks Using Complex Network Techniques and Natural Language Processing}\\[0.35em]
{\itshape Mohammad Farhadishad, Mohammad Khansari}\\[0.2em]
{\footnotesize First author: Mohammad Farhadishad}\\[0.45em]
\noindent\rule{\textwidth}{0.3pt}\\[0.4em]
{\footnotesize  Many real-world systems can be effectively modeled and analyzed using graph-based approaches. Despite the growing interest in legal network analysis, studies on Persian legal networks remain limited, particularly those integrating complex network science and artificial intelligence. This study investigates a large-scale Persian legal network constructed from 4,921 judicial rulings, 15 charge titles, and 807 legal documents from the criminal courts of the Islamic Republic of Iran. While the original dataset contained only judicial rulings, charge titles and legal documents were automatically extracted from ruling texts using natural language processing and an AI-based question-answering model. The extracted entities were used to construct and analyze bipartite and tripartite legal networks. Results show that the network exhibits several structural properties commonly observed in real-world complex networks. Furthermore, the proposed multi-level representation reveals meaningful relationships among judicial rulings, charge titles, and legal documents, providing valuable insights for legal network analysis.}\\[0.4em]
{\footnotesize\bfseries Keywords:} {\footnotesize Complex networks, multipartite graphs, legal networks, natural language processing, intelligent question and answer. \textbackslash\{\}\textbackslash\{\}}\\[0.15em]
{\footnotesize\bfseries MSC:} {\footnotesize 99X99 , 99X99 , 99X99 . 1cm}\\
\vfill
\end{minipage}
\newpage

% p.1331
\noindent\begin{minipage}[t][\textheight]{\textwidth}
\noindent{\bfseries\large Information Loss in Spatial Statistics Analysis: Entropy, Scale, and Uncertainty}\\[0.35em]
{\itshape Zahra Amini Farsani}\\[0.2em]
{\footnotesize First author: Zahra Amini Farsani}\\[0.45em]
\noindent\rule{\textwidth}{0.3pt}\\[0.4em]
{\footnotesize  Spatial data often contain several layers of information that are not fully represented in final statistical summaries. A spatial field, image, disease map, or set of regional observations may include spatial dependence, local heterogeneity, scale effects, and uncertainty. In statistical practice, however, these complex structures are often reduced to estimated parameters, smoothed surfaces, clusters, risk measures, test results, or prediction errors. Such reductions are useful and sometimes necessary for interpretation, communication, and decision-making, but they may also hide important parts of the original spatial information. This talk discusses information loss in spatial statistics from an entropy-based perspective. Entropy provides a natural way to think about information, uncertainty, and reduction in spatially structured data. The aim is to examine how spatial information changes when data are aggregated, smoothed, classified, or summarized, and to highlight the role of entropy, Bayesian inference, and uncertainty quantification in understanding what is preserved, simplified, or lost during spatial statistical analysis.}\\[0.4em]
{\footnotesize\bfseries Keywords:} {\footnotesize Spatial statistics, Entropy, Information loss, Spatial dependence, Scale effects, Uncertainty quantification. \textbackslash\{\}}\\[0.15em]
{\footnotesize\bfseries MSC:} {\footnotesize 62M30 , 94A17 , 62F15 . \textbackslash\{\} 1cm}\\
\vfill
\end{minipage}
\newpage

% p.1200
\noindent\begin{minipage}[t][\textheight]{\textwidth}
\noindent{\bfseries\large Comparative Analysis of Bayesian Estimators for the MOAP-IL Distribution under Different Loss Functions}\\[0.35em]
{\itshape N.Sanjari Farsipour}\\[0.2em]
{\footnotesize First author: N.Sanjari Farsipour}\\[0.45em]
\noindent\rule{\textwidth}{0.3pt}\\[0.4em]
{\footnotesize  The Marshall--Olkin Alpha Power Inverse Lindley (MOAP-IL) distribution is a flexible model for lifetime data with non-monotonic hazard rates. In this paper, we propose a Bayesian approach to jointly estimate the shape parameter and the rate parameter . Because the posterior distribution is nonlinear and has no closed-form solution, we compute Bayesian estimates using MCMC with the Metropolis--Hastings algorithm. We study 13 loss functions in three groups: symmetric, asymmetric, and relative. Simulation results show that symmetric loss functions (such as SEL and Stein) usually give smaller MSE, while asymmetric loss functions (such as AREL and M-Linex) tend to produce more conservative estimates. This difference is important for risk-sensitive settings where overestimation and underestimation do not have the same cost. Convergence diagnostics, including trace and density plots, suggest stable estimation and good mixing, which supports the reliability of the proposed method.}\\[0.4em]
{\footnotesize\bfseries Keywords:} {\footnotesize Bayesian inference, MOAP-IL distribution, Asymmetric loss functions, MCMC algorithm, and Markov chain convergence \textbackslash\{\}\textbackslash\{\}}\\[0.15em]
{\footnotesize\bfseries MSC:} {\footnotesize 62F15 , 62F10 . 1cm}\\
\vfill
\end{minipage}
\newpage

% p.1119
\noindent\begin{minipage}[t][\textheight]{\textwidth}
\noindent{\bfseries\large Bayesian Inference and Tracking of System Reliability}\\[0.35em]
{\itshape Masoud Ganji}\\[0.2em]
{\footnotesize First author: Masoud Ganji}\\[0.45em]
\noindent\rule{\textwidth}{0.3pt}\\[0.4em]
{\footnotesize  This paper develops an extended Bayesian framework for estimating and tracking system reliability with a clear focus on sequential use of data and prior information. Within this framework, we establish new theoretical results, including the uniqueness of Beta priors under moment matching, bounds for series--parallel system reliability, asymptotic normality of posterior distributions, and Bayesian consistency. We further derive limit theorems for posterior mean convergence rates and credible-interval bounds for mixed systems. Numerical examples and sequential updating schemes demonstrate how the proposed framework can be applied in practice. Overall, the approach provides a rigorous, computationally light, and reproducible methodology for reliability analysis.}\\[0.4em]
{\footnotesize\bfseries Keywords:} {\footnotesize Bayesian Inference, System Reliability, Posterior Consistency, Credible Bounds}\\[0.15em]
{\footnotesize\bfseries MSC:} {\footnotesize 62F15 , 62N05 , 62P30 , 93B30 .}\\
\vfill
\end{minipage}
\newpage

% p.1118
\noindent\begin{minipage}[t][\textheight]{\textwidth}
\noindent{\bfseries\large Bayesian Parameter and Reliability Estimation for a Bivariate Weibull Distribution: Parallel System}\\[0.35em]
{\itshape Masoud Ganji}\\[0.2em]
{\footnotesize First author: Masoud Ganji}\\[0.45em]
\noindent\rule{\textwidth}{0.3pt}\\[0.4em]
{\footnotesize  This study develops a reliability modeling framework for multi-component systems using a Bivariate Weibull distribution under both independence and dependence assumptions. Dependence is introduced through the Gumbel copula, and parameters are estimated via maximum likelihood and Bayesian methods. Simulations show that independence-based models perform well only when components are truly independent, whereas ignoring dependence leads to biased reliability estimates. The copula-based model more accurately captures dependence, and Bayesian inference provides estimates comparable to maximum likelihood while offering credible intervals for parameters and system reliability. \textbackslash\{\}\textbackslash\{\}}\\[0.4em]
{\footnotesize\bfseries Keywords:} {\footnotesize Bivariate Weibull, Bayes estimation, Parallel system reliability, Copula models, MCMC \textbackslash\{\}\textbackslash\{\}}\\[0.15em]
{\footnotesize\bfseries MSC:} {\footnotesize 62F15 , 62N05 , 62P30 , 93B30 .}\\
\vfill
\end{minipage}
\newpage

% p.1319
\noindent\begin{minipage}[t][\textheight]{\textwidth}
\noindent{\bfseries\large Optimal plan under Type-II Censoring Schemes in Multi-component Stress-Strength Reliability based on its Estimate and Cost}\\[0.35em]
{\itshape Somayeh Ghafouri, Elham Basiri}\\[0.2em]
{\footnotesize First author: Somayeh Ghafouri}\\[0.45em]
\noindent\rule{\textwidth}{0.3pt}\\[0.4em]
{\footnotesize  In this paper, we consider a cut point M\textasciicircum{}* for number of failure times in the stress component of a multi-component stress-strength system based on progressively Type-II censoring assumed as a truncated geometry distribution while the number of failure times in strength component is assumed fixed. The optimal schemes and the cut points can be considered and the most important is determination them so that the total cost of the test and the estimation of multi-component stress-strength reliability have been reached to minimum and maximum, respectively. A simulation study is carried out to check which different censoring schemes and cut points have the best performance in terms of the total cost and the estimation of multi-component stress-strength reliability. At the end, a real example is given to illustrate the proposed method.\textbackslash\{\}\textbackslash\{\}}\\[0.4em]
{\footnotesize\bfseries Keywords:} {\footnotesize Multi-component stress-strength reliability, Type-II right censoring, Truncated geometry distribution. \textbackslash\{\}\textbackslash\{\}}\\[0.15em]
{\footnotesize\bfseries MSC:} {\footnotesize 62N01 , 62N05 . 1cm}\\
\vfill
\end{minipage}
\newpage

% p.1083
\noindent\begin{minipage}[t][\textheight]{\textwidth}
\noindent{\bfseries\large series Explainable Ensemble Learning for Software Defect Prediction [4pt] Using Fuzzy Integral-Based Aggregation}\\[0.35em]
{\itshape Parviz Ghorbanzadeh}\\[0.2em]
{\footnotesize First author: Parviz Ghorbanzadeh}\\[0.45em]
\noindent\rule{\textwidth}{0.3pt}\\[0.4em]
{\footnotesize The adoption of artificial intelligence in software defect prediction (SDP) has significantly improved predictive performance; however, the lack of transparency in complex models limits their practical applicability in critical software engineering environments. This paper proposes an explainable ensemble learning framework that integrates multiple base classifiers using a fuzzy integral-based aggregation mechanism. By modelling the interaction among classifiers through fuzzy measures, the approach captures both individual and collective predictive contributions. Additionally, interpretability is achieved by deriving importance indices that quantify the influence of features and models on final decisions. Experiments conducted on five NASA benchmark datasets (CM1, JM1, KC1, KC2, PC1) demonstrate that the proposed method achieves superior performance in terms of AUC, F-measure, and Matthews Correlation Coefficient (MCC) while providing meaningful explanations for prediction outcomes. The results highlight the potential of fuzzy-integral-based explainable AI in enhancing trust and reliability in SDP}\\[0.4em]
{\footnotesize\bfseries Keywords:} {\footnotesize Software Defect Prediction; Explainable AI; Fuzzy Integral; Ensemble Learning; Choquet Integral. 6pt 10pt @twocolumnfalse ]}\\[0.15em]
\vfill
\end{minipage}
\newpage

% p.1151
\noindent\begin{minipage}[t][\textheight]{\textwidth}
\noindent{\bfseries\large Deep Learning Helping to Realize Masnavi-e Maanavi}\\[0.35em]
{\itshape Mousa Golalizadeh}\\[0.2em]
{\footnotesize First author: Mousa Golalizadeh}\\[0.45em]
\noindent\rule{\textwidth}{0.3pt}\\[0.4em]
{\footnotesize  This study introduces a novel quantitative methodology, borrowed from deep learning, to validate the hypothesis that Neynameh serves as a condensed summary of Masnavi-e Maanavi. We employ a multi-layered analytical framework combining ParsBERT fine-tuned on classical Persian texts with semantic embeddings, distribution analysis, spiritual journey kapping, and thematic analysis. Based on 18 verses of Neynameh and 26000 verses from Masnavi-e Maanavi, our analysis revealed a mean similarity score of 0.774 with relatively uniform distribution. The composite structural analysis yielded a score of 0.785, significantly above the validation threshold of 0.75. This methodology offers a replicable framework applicable to other texts in Persian literary tradition.}\\[0.4em]
{\footnotesize\bfseries Keywords:} {\footnotesize Deep learning, Multi-layered, Masnavi-e Maanavi, ParsBERT, Similarity score. \textbackslash\{\}\textbackslash\{\}}\\[0.15em]
{\footnotesize\bfseries MSC:} {\footnotesize 62P99 , 62P25 . 1cm}\\
\vfill
\end{minipage}
\newpage

% p.1245
\noindent\begin{minipage}[t][\textheight]{\textwidth}
\noindent{\bfseries\large Ordered Ranked Set Sampling for Progressive-Stress Accelerated Life Tests under Type-II Censoring}\\[0.35em]
{\itshape Nooshin Hakamipour}\\[0.2em]
{\footnotesize First author: Nooshin Hakamipour}\\[0.45em]
\noindent\rule{\textwidth}{0.3pt}\\[0.4em]
{\footnotesize  This paper investigates the performance of the ordered ranked set sampling scheme under progressive-stress accelerated life tests with type-II censoring. The lifetime of the item under use stress is assumed to follow a Weibull distribution, while the scale parameter is related to the applied stress through the inverse power law. Maximum likelihood estimators of the model parameters are derived under both ORSS and simple random sampling schemes. The finite-sample behavior of the proposed estimators is examined through simulation studies and numerical results. The findings indicate that ORSS yields more efficient estimation than SRS in terms of smaller absolute biases and mean squared errors.}\\[0.4em]
{\footnotesize\bfseries Keywords:} {\footnotesize Maximum Likelihood Estimation, Ordered Ranked Set Sampling, Progressive-Stress Accelerated Life Test, Type-II Censoring. \textbackslash\{\}\textbackslash\{\}}\\[0.15em]
{\footnotesize\bfseries MSC:} {\footnotesize 99X99 , 99X99 , 99X99 . 1cm}\\
\vfill
\end{minipage}
\newpage

% p.1186
\noindent\begin{minipage}[t][\textheight]{\textwidth}
\noindent{\bfseries\large Mixture modeling of normal mean-variance Birnbaum-Saunders factor analyzers under a Bayesian framework}\\[0.35em]
{\itshape Farzane Hashemi, Fatemeh Karimi}\\[0.2em]
{\footnotesize First author: Farzane Hashemi}\\[0.45em]
\noindent\rule{\textwidth}{0.3pt}\\[0.4em]
{\footnotesize  The mixture of factor analyzers (MFA) model, by representing component covariance matrices through a factor-analytic structure, efficiently reduces the number of free parameters and is widely used to uncover latent groups in high-dimensional data. While extensions to handle skewed data and asymmetry within latent groups have been explored in frequentest frameworks, they often face computational limitations. Here, we propose a Bayesian approach for normal mean-variance Birnbaum-Saunders MFA models, demonstrating its flexibility and performance on real dataset and highlighting computational advantages.}\\[0.4em]
{\footnotesize\bfseries Keywords:} {\footnotesize Bayesian analysis, Mixture of factor analysis model, Normal mean-variance Birnbaum-Saunders distribution, and Seeds data. \textbackslash\{\}\textbackslash\{\}}\\[0.15em]
{\footnotesize\bfseries MSC:} {\footnotesize 99X99 , 99X99 , 99X99 . \textbackslash\{\}\textbackslash\{\} 0.5cm}\\
\vfill
\end{minipage}
\newpage

% p.1230
\noindent\begin{minipage}[t][\textheight]{\textwidth}
\noindent{\bfseries\large Stress-strength reliability analysis under partially accelerated life testing using the inverse Weibull distribution}\\[0.35em]
{\itshape Farzane Hashemi}\\[0.2em]
{\footnotesize First author: Farzane Hashemi}\\[0.45em]
\noindent\rule{\textwidth}{0.3pt}\\[0.4em]
{\footnotesize  The stress-strength reliability of a system is defined as the probability that its strength exceeds the applied stress, namely R=P(X>Y) . In this paper, we study the estimation of R when the strength X and stress Y are independent random variables following inverse Weibull distributions, and the strength variable is subjected to a step-stress partially accelerated life test. The inverse Weibull model provides a flexible alternative to lifetime distributions with restrictive failure-rate behavior. Since the maximum likelihood estimator of R and its asymptotic distribution cannot be obtained in closed form, numerical estimation and asymptotic confidence intervals are considered. A real dataset is analyzed to illustrate the applicability and performance of the proposed model.}\\[0.4em]
{\footnotesize\bfseries Keywords:} {\footnotesize Partial acceleration, Inverse Weibull distributions, ML method, and Reliability analysis. \textbackslash\{\}\textbackslash\{\}}\\[0.15em]
{\footnotesize\bfseries MSC:} {\footnotesize 99X99 , 99X99 , 99X99 . \textbackslash\{\}\textbackslash\{\} 0.1cm}\\
\vfill
\end{minipage}
\newpage

% p.1050
\noindent\begin{minipage}[t][\textheight]{\textwidth}
\noindent{\bfseries\large Performance of the auto.arima( Function in Recovering True ARIMA Models Order Under Three Different Stationarity Tests}\\[0.35em]
{\itshape M. Heidari}\\[0.2em]
{\footnotesize First author: M. Heidari}\\[0.45em]
\noindent\rule{\textwidth}{0.3pt}\\[0.4em]
{\footnotesize The auto.arima() function in R is widely used for automated ARIMA model selection in time series analysis. However, its performance may depending on the choice of stationarity test for determining the order ( p,d,q ). This study investigates how the auto.arima() function behaves under three unit root tests (three stationarity tests); Augmented Dickey-Fuller (ADF), Kwiatkowski--Phillips--Schmidt--Shin (KPSS), and Phillips--Perron (PP), when determining the correct ARIMA model order (p, d, q) . Using 10 simulated ARIMA (p,d,q) models with varying parameters, we compute the percentage of correct identification of (p,d,q) simultaneously for each test. Among the three unit root tests, auto.arima() function on KPSS achieves the highest mean correct identification rate ( 29 -- 35\textbackslash\{\}}\\[0.4em]
{\footnotesize\bfseries Keywords:} {\footnotesize ARIMA model, auto.arima(), ADF test, KPSS test, PP test, time series and stationarity.\textbackslash\{\}\textbackslash\{\}}\\[0.15em]
{\footnotesize\bfseries MSC:} {\footnotesize 62M10 , 62F03 . 1cm}\\
\vfill
\end{minipage}
\newpage

% p.1232
\noindent\begin{minipage}[t][\textheight]{\textwidth}
\noindent{\bfseries\large Generalized H}\\[0.35em]
{\itshape Hazhir Homei, Alireza Hedayati, Maryam Kamali Alamdari}\\[0.2em]
{\footnotesize First author: Hazhir Homei}\\[0.45em]
\noindent\rule{\textwidth}{0.3pt}\\[0.4em]
{\footnotesize  This paper introduces the concept of an H -sufficient statistic and explores its theoretical implications in the derivation of uniformly minimum variance unbiased estimators (UMVUEs). By leveraging this framework, we extend classical results associated with the Rao-Blackwell and Lehmann-Scheffé theorems, particularly in settings where a complete sufficient statistic is unavailable. The proposed approach not only clarifies several longstanding examples in statistical literature but also yields novel interpretations of optimality conditions. Furthermore, we establish connections between H -sufficiency and Basu's theorem, demonstrating how independence and uncorrelatedness criteria can be systematically utilized to identify optimal estimators. The results provide a unified and flexible methodology for handling incomplete minimal sufficient statistics in parametric inference.}\\[0.4em]
{\footnotesize\bfseries Keywords:} {\footnotesize Rao-Blackwell theorem, Lehmann-Scheffé theorem, H -sufficient statistic, incomplete sufficient statistics, Basu's theorem, UMVUE. \textbackslash\{\}\textbackslash\{\}}\\[0.15em]
{\footnotesize\bfseries MSC:} {\footnotesize 62F10 , 62B05 , 62F99 . 1cm}\\
\vfill
\end{minipage}
\newpage

% p.1022
\noindent\begin{minipage}[t][\textheight]{\textwidth}
\noindent{\bfseries\large A Comparative Performance Analysis of Structural Equation Modeling Approaches: A Monte Carlo Simulation Study}\\[0.35em]
{\itshape Jamshid Jamali}\\[0.2em]
{\footnotesize First author: Jamshid Jamali}\\[0.45em]
\noindent\rule{\textwidth}{0.3pt}\\[0.4em]
{\footnotesize  This Monte Carlo simulation study compared the performance of three structural equation modeling approaches CB-SEM, PLS-SEM, and GSCA within a mediation model under varying sample sizes, data distributions, and indicator counts. Using full factorial design with 500 replications per scenario, results revealed that CB-SEM provided the most accurate estimates under normality and larger samples \textbackslash\{\}( ( N 250 ) \textbackslash\{\}), yet suffered convergence issues in small or severely non-normal conditions. PLS-SEM demonstrated superior convergence rates and statistical power across all scenarios, albeit with systematic attenuation bias. GSCA offered balanced intermediate performance, particularly robust under distributional violations. Increasing indicators per construct benefited CB-SEM most substantially. No single method proved universally superior; rather, selection should align with research objectives, sample availability, and data characteristics. These evidence-based findings offer practical guidance for biostatisticians and health researchers in choosing appropriate SEM methodologies.}\\[0.4em]
{\footnotesize\bfseries Keywords:} {\footnotesize CB-SEM, PLS-SEM, GSCA, Monte Carlo Simulation,Mediation Analysis. \textbackslash\{\}\textbackslash\{\}}\\[0.15em]
{\footnotesize\bfseries MSC:} {\footnotesize 62P10. 1cm}\\
\vfill
\end{minipage}
\newpage

% p.1131
\noindent\begin{minipage}[t][\textheight]{\textwidth}
\noindent{\bfseries\large Estimating the Stress-Strength Parameter Using a Dirichlet Process Mixture Model with Burr XII Kernels}\\[0.35em]
{\itshape Bahram Haji Joudaki}\\[0.2em]
{\footnotesize First author: Bahram Haji Joudaki}\\[0.45em]
\noindent\rule{\textwidth}{0.3pt}\\[0.4em]
{\footnotesize  In this article, a new model is presented for modeling the distribution of stress and strength variables in a system. In this model, it is assumed that the distributions of the stress and strength variables are independent of one another, and a mixture of Dirichlet processes with Burr XII kernels is used to model the distribution of each of these variables. Prior distributions are selected for the parameters present in these models. The posterior distributions of the parameters have a complex form; therefore, Gibbs sampling is employed to generate samples from these distributions. The Dirichletprocess package in R software is used to implement Gibbs algorithms in the mentioned mixture models. Under this new model, the stress-strength parameter is estimated. To evaluate the performance of the proposed model, a real datasets is utilized. The results obtained demonstrate satisfactory performance of the proposed model.}\\[0.4em]
{\footnotesize\bfseries Keywords:} {\footnotesize Burr XII, Dirichlet process, Gibbs algorithm, Mixture model, Stress-Strength}\\[0.15em]
{\footnotesize\bfseries MSC:} {\footnotesize 62F15 , 62G05 . 1cm}\\
\vfill
\end{minipage}
\newpage

% p.1085
\noindent\begin{minipage}[t][\textheight]{\textwidth}
\noindent{\bfseries\large Shrinkage Bayes Estimation for Matrix-Variate Normal Mean Matrix under Balanced Loss Function}\\[0.35em]
{\itshape Hamid Karamikabir h}\\[0.2em]
{\footnotesize First author: Hamid Karamikabir h}\\[0.45em]
\noindent\rule{\textwidth}{0.3pt}\\[0.4em]
{\footnotesize  We study Bayesian estimation of the mean matrix in a matrix-variate normal distribution with known covariances and a conjugate prior. Under a balanced matrix composite loss function, we derive the Bayes estimator as a convex combination of the posterior mean and the observed data, controlled by a shrinkage parameter [0,1] . The posterior mean formula \_n = 1 2 ( \_0 + X) is established. A simulation study using non-diagonal AR(1) covariances confirms the theoretical findings, showing that the optimal lies strictly between zero and one under high noise, with the optimal estimator achieving 38.4\textbackslash\{\}}\\[0.4em]
{\footnotesize\bfseries Keywords:} {\footnotesize Balanced loss function, Bayes estimator, Image denoising, Matrix-variate normal distribution, Shrinkage estimation. \textbackslash\{\}\textbackslash\{\}}\\[0.15em]
{\footnotesize\bfseries MSC:} {\footnotesize 62J07 , 62F15 . 1cm}\\
\vfill
\end{minipage}
\newpage

% p.1128
\noindent\begin{minipage}[t][\textheight]{\textwidth}
\noindent{\bfseries\large Efficient Bayesian Inference for Stress-Strength Reliability in Gompertz Distribution via Ranked Set Sampling}\\[0.35em]
{\itshape Z. Karimiezmareh}\\[0.2em]
{\footnotesize First author: Z. Karimiezmareh}\\[0.45em]
\noindent\rule{\textwidth}{0.3pt}\\[0.4em]
{\footnotesize  The Gompertz distribution, with its exponentially increasing hazard rate, is widely used in reliability engineering and survival analysis. However, estimating stress-strength reliability under this distribution has received limited attention. Conventional simple random sampling (SRS) is inefficient when precise measurements are costly or time-consuming. Ranked set sampling (RSS) addresses this limitation by using inexpensive ranking before measurement, extracting more information without increasing sample size. This study develops a Bayesian estimation of stress-strength reliability R=P(X<Y) for the Gompertz distribution under RSS. Bayesian estimators are derived under a squared error loss function using the Metropolis-Hastings algorithm. The significance of this work is providing an efficient, low-cost, and high-precision method for reliability estimation when accurate measurement is expensive or difficult.}\\[0.4em]
{\footnotesize\bfseries Keywords:} {\footnotesize Ranked Set Sampling, Gompertz Distribution, Stress-Strength Reliability, and Metropolics-Hastings Algorithm. \textbackslash\{\}\textbackslash\{\}}\\[0.15em]
{\footnotesize\bfseries MSC:} {\footnotesize 62Exx , 62F10 , 62F15 . 1cm}\\
\vfill
\end{minipage}
\newpage

% p.1129
\noindent\begin{minipage}[t][\textheight]{\textwidth}
\noindent{\bfseries\large A Machine Learning-Driven Framework for the Implementation of Statistical Population Registers: A Hybrid Intelligent Approach}\\[0.35em]
{\itshape Vadood Keramati, Sheida Abdolmaleki}\\[0.2em]
{\footnotesize First author: Vadood Keramati}\\[0.45em]
\noindent\rule{\textwidth}{0.3pt}\\[0.4em]
{\footnotesize  The rapid evolution of data ecosystems has necessitated a paradigm shift in how national statistical organizations construct and maintain population registers. The traditional reliance on costly, periodic censuses is increasingly being replaced by the implementation of Statistical Population Registers (SPR) derived from multi-source administrative data. However, the operationalization of an SPR is not merely a data aggregation task; it is a complex engineering challenge requiring the reconciliation of heterogeneous datasets characterized by noise, missing identifiers, and structural variations. In the context of Persian administrative registries, these implementation challenges are particularly pronounced. The lack of standardized orthography and the absence of a unified, error-free identifier create significant barriers to effective record integration. Conventional deterministic methods often prove inadequate in handling such linguistic and structural complexities, while purely probabilistic linkage techniques frequently suffer from computational bottlenecks and suboptimal classification precision. This research addresses these challenges by proposing a robust, machine learning-driven framework for the implementation of an SPR. By introducing a hybrid intelligent approach that synergizes the structural rigor of Probabilistic Record Linkage (PRL) with the adaptive predictive power of supervised machine learning---specifically decision tree classifiers---this study provides a scalable solution for high-fidelity data integration. The central objective of this research is to establish an automated pipeline that minimizes human-in-the-loop requirements while maximizing linkage accuracy. This paper evaluates the proposed framework through rigorous empirical testing on simulated administrative datasets, demonstrating how the integration of machine learning leads to superior performance in the implementation of reliable, modern statistical registers.}\\[0.4em]
{\footnotesize\bfseries Keywords:} {\footnotesize Intelligent System, Machine Learning, Administrative Data Integration, Record Linkage, Statistical Population Register, Persian Language Processing. \textbackslash\{\}\textbackslash\{\}}\\[0.15em]
{\footnotesize\bfseries MSC:} {\footnotesize 68T05 , 62P25 , 68U35 . 1cm}\\
\vfill
\end{minipage}
\newpage

% p.1295
\noindent\begin{minipage}[t][\textheight]{\textwidth}
\noindent{\bfseries\large The Asymmetric Laplace Process: A Powerful Alternative to AL in Spatial Quantile Regression}\\[0.35em]
{\itshape Sepehr Khademhosseini, Firoozeh Rivaz s}\\[0.2em]
{\footnotesize First author: Sepehr Khademhosseini}\\[0.45em]
\noindent\rule{\textwidth}{0.3pt}\\[0.4em]
{\footnotesize  Quantile regression is a powerful tool for modelling the conditional distribution of a response variable, particularly when interest lies in the tails of the distribution. Standard quantile regression models typically employ the asymmetric Laplace (AL) distribution under the assumption of independent errors, and therefore cannot capture spatial dependence when it exists in the data. To address this, Lum and Gelfand (2012) introduced the Asymmetric Laplace Process (ALP), which incorporates a latent Gaussian process to capture spatial correlation within the error structure. We conduct a simulation-based comparison evaluating the ALP against the standard AL model under varying levels of spatial correlation. Results show that the ALP achieves lower RMSE for the regression parameter in almost all scenarios. However, credible interval coverage for \_0 declines substantially as the spatial range increases, ultimately falling below the nominal level. This deterioration in coverage suggests a potential identifiability issue between the spatial process and the intercept under strong spatial correlation. We conclude that while ALP offers improved point estimation, its credible intervals should be interpreted with caution in the presence of strong spatial dependence.}\\[0.4em]
{\footnotesize\bfseries Keywords:} {\footnotesize Spatial quantile regression, Asymmetric Laplace process, Bayesian inference, Stan.\textbackslash\{\}\textbackslash\{\}}\\[0.15em]
{\footnotesize\bfseries MSC:} {\footnotesize 62F15 , 62M30 , 62J05 . 1cm}\\
\vfill
\end{minipage}
\newpage

% p.1258
\noindent\begin{minipage}[t][\textheight]{\textwidth}
\noindent{\bfseries\large Evaluating Large Language Models on the Iranian University Entrance Exam: A Multi-Subject Persian Benchmark}\\[0.35em]
{\itshape Shayan Kheirabadi}\\[0.2em]
{\footnotesize First author: Shayan Kheirabadi}\\[0.45em]
\noindent\rule{\textwidth}{0.3pt}\\[0.4em]
{\footnotesize  This paper presents a comprehensive benchmark of state-of-the-art Large Language Models (LLMs)---specifically Gemini v3, GPT-5.1, and DeepSeek 3.2---on the 2025 Iranian National Entrance Exam (Konkur). We evaluated the zero-shot performance of these models across three major academic streams: Mathematics, Experimental Sciences, and Humanities. The results reveal a significant "Specialization Gap." GPT-5.1 demonstrated superior logical reasoning capabilities, achieving the highest accuracy in Mathematics (50.0\textbackslash\{\} 0.1cm}\\[0.4em]
{\footnotesize\bfseries Keywords:} {\footnotesize Large Language Models (LLMs , Persian Benchmark, Iranian National Entrance Exam (Konkur , GPT-5.1, Gemini v3, DeepSeek 3.2, Zero-shot Learning, Multimodal Evaluation. 1cm}\\[0.15em]
\vfill
\end{minipage}
\newpage

% p.1098
\noindent\begin{minipage}[t][\textheight]{\textwidth}
\noindent{\bfseries\large Optimal Designs under a Proportional Interference Model with Nonnegative Proportionality Constant and Correlated Observations}\\[0.35em]
{\itshape Razieh Khodsiani}\\[0.2em]
{\footnotesize First author: Razieh Khodsiani}\\[0.45em]
\noindent\rule{\textwidth}{0.3pt}\\[0.4em]
{\footnotesize  In block experiments where a treatment applied to one plot affects neighboring plots, interference models are used to account for these neighbor effects. A special class is the proportional interference model, in which each neighbor effect is a proportion of the direct treatment effect. We study a proportional one-sided interference model in which only the left neighbor contributes a positive proportion of the direct treatment effect, and observations follow a first-order circular autoregressive correlation structure. Under this model we derive exact optimal designs for correlated observations.}\\[0.4em]
{\footnotesize\bfseries Keywords:} {\footnotesize Interference models, Circular block designs, Neighbor balanced designs, Universally optimal criterion.\textbackslash\{\}\textbackslash\{\}}\\[0.15em]
{\footnotesize\bfseries MSC:} {\footnotesize 62K05 , 62K10 . 1cm}\\
\vfill
\end{minipage}
\newpage

% p.1178
\noindent\begin{minipage}[t][\textheight]{\textwidth}
\noindent{\bfseries\large Estimation of the Mean Residual Life Function Using Neoteric Ranked Set Sampling}\\[0.35em]
{\itshape Leila Jabari Koopaei}\\[0.2em]
{\footnotesize First author: Leila Jabari Koopaei}\\[0.45em]
\noindent\rule{\textwidth}{0.3pt}\\[0.4em]
{\footnotesize  The mean residual life (MRL) function is one of the fundamental concepts in survival analysis. The MRL is easily interpretable and particularly useful for conveying the survival information to patients. Ranked set sampling (RSS) is a sampling design useful for situations where precise quantification of units is costly or difficult, but ranking them based on associated variables is easy with negligible cost. In recent years, different modifications of RSS have been introduced to enhance its efficiency. This paper investigates the estimation of the MRL function using Neoteric Ranked Set Sampling (NRSS), a notable modification of RSS. We compare the proposed estimator with its counterparts in RSS and SRS using Monte Carlo simulation. Our results show that the proposed estimator is more efficient than its counterparts for small and moderate sample sizes.}\\[0.4em]
{\footnotesize\bfseries Keywords:} {\footnotesize Mean residual life, Neoteric ranked set sampling, Monte Carlo simulation, Ranked set sampling. \textbackslash\{\}\textbackslash\{\}}\\[0.15em]
{\footnotesize\bfseries MSC:} {\footnotesize 62D05 , 62G05 . 1cm}\\
\vfill
\end{minipage}
\newpage

% p.1347
\noindent\begin{minipage}[t][\textheight]{\textwidth}
\noindent{\bfseries\large A Review of Smooth Estimation of the Cumulative Distribution Function under Judgment Post Stratification}\\[0.35em]
{\itshape Azizi Kouhanestani, M., Zamanzade, E., Goli, S.}\\[0.2em]
{\footnotesize First author: Azizi Kouhanestani}\\[0.45em]
\noindent\rule{\textwidth}{0.3pt}\\[0.4em]
{\footnotesize  This paper provides a review of smooth estimation methods for the cumulative distribution function (CDF) under the judgment post stratification (JPS) sampling scheme. Particular attention is devoted to a general class of CDF estimators that encompasses both empirical and kernel-based estimators. The main finite-sample and asymptotic properties of these estimators are summarized and discussed. To further illustrate their performance, an extended Monte Carlo simulation study is conducted to compare the kernel distribution function (KDF) estimator under JPS with its counterpart based on simple random sampling (SRS). The comparison considers a range of sample sizes, set sizes, and ranking qualities. The findings demonstrate that incorporating ranking information through the JPS design can lead to considerable efficiency gains relative to SRS across a broad spectrum of settings.}\\[0.4em]
{\footnotesize\bfseries Keywords:} {\footnotesize Nonparametric estimation, Monte Carlo simulation, Judgment ranking, and Statistical inference \textbackslash\{\}\textbackslash\{\}}\\[0.15em]
{\footnotesize\bfseries MSC:} {\footnotesize 62D05 , 62G05 . 1cm}\\
\vfill
\end{minipage}
\newpage

% p.1299
\noindent\begin{minipage}[t][\textheight]{\textwidth}
\noindent{\bfseries\large Preliminaries and Problem Formulation}\\[0.35em]
{\itshape Javad Koushki}\\[0.2em]
{\footnotesize First author: Javad Koushki}\\[0.45em]
\noindent\rule{\textwidth}{0.3pt}\\[0.4em]
{\footnotesize  Proper efficiency is a foundational solution concept in deterministic multiobjective optimization. Recently, the notion of anchor points was introduced to classify properly efficient solutions that solve a weighted sum scalarization problem where at least one objective weight is zero. In practical applications, however, objectives are frequently corrupted by stochastic uncertainty. This paper generalizes the concept of anchor points to the framework of Stochastic Multiobjective Optimization (SMOP) using a probabilistic expected value approach. We introduce the concept of Stochastic Anchor Points and characterize them via properties of the stochastic objective space, the stochastic tangent cone, and the stochastic normal cone.}\\[0.4em]
{\footnotesize\bfseries Keywords:} {\footnotesize Stochastic multiobjective optimization; Anchor points; Proper efficiency; Expected value scalarization; Normal cone; Uncertainty analysis. \textbackslash\{\}\textbackslash\{\}}\\[0.15em]
{\footnotesize\bfseries MSC:} {\footnotesize 99X99 , 99X99 , 99X99 . \textbackslash\{\}\textbackslash\{\}}\\
\vfill
\end{minipage}
\newpage

% p.1004
\noindent\begin{minipage}[t][\textheight]{\textwidth}
\noindent{\bfseries\large Investigating the Relationship between Iranian Insurance Companies' Solvency Ratio and Insurance Policy Sales}\\[0.35em]
{\itshape Rahim Mahmoudvand}\\[0.2em]
{\footnotesize First author: Rahim Mahmoudvand}\\[0.45em]
\noindent\rule{\textwidth}{0.3pt}\\[0.4em]
{\footnotesize  The Central Insurance of Iran introduced the financial solvency ratio as a supervisory tool for insurance companies following the approval of Regulation No. 69 in February 2012. This metric assesses an insurer's financial capacity to cover accepted risks and classifies companies into five levels, with Level 1 (ratio above 100\textbackslash\{\}}\\[0.4em]
{\footnotesize\bfseries Keywords:} {\footnotesize Solvency Ratio, Insurance Regulation, Market Discipline, and Policy \textbackslash\{\}\textbackslash\{\}}\\[0.15em]
{\footnotesize\bfseries MSC:} {\footnotesize 62P05 , 91B30 , 62H20 . 1cm}\\
\vfill
\end{minipage}
\newpage

% p.1339
\noindent\begin{minipage}[t][\textheight]{\textwidth}
\noindent{\bfseries\large Bayesian Inference for VAR-GoGARCH Models under Gaussian and Skew-Normal Innovations}\\[0.35em]
{\itshape T. Manouchehri , A. R. Nematollahi, M. Sadeghi Kahmini}\\[0.2em]
{\footnotesize First author: T. Manouchehri}\\[0.45em]
\noindent\rule{\textwidth}{0.3pt}\\[0.4em]
{\footnotesize  Modeling multivariate financial time series requires flexible frameworks capable of capturing dynamic dependence, conditional heteroscedasticity, and asymmetric innovation distributions. This paper develops a Bayesian estimation framework for VAR-GoGARCH models under both Gaussian and skew-normal innovations. Posterior inference is performed using a Metropolis--Hastings Markov chain Monte Carlo (MCMC) algorithm, while the mixing matrix is estimated separately via the SOBI decomposition. The proposed methodology is evaluated through Monte Carlo experiments under correctly specified and asymmetric innovation settings. The results show that the Bayesian estimator consistently achieves lower root mean squared errors than the classical maximum likelihood estimator. Moreover, when the innovation distribution is correctly specified as skew-normal, the proposed approach substantially improves volatility estimation, reducing the mean squared error by approximately 59\textbackslash\{\}}\\[0.4em]
{\footnotesize\bfseries Keywords:} {\footnotesize Bayesian inference, VAR-GoGARCH model, Skew-normal distribution, Markov chain Monte Carlo, Volatility estimation. \textbackslash\{\}\textbackslash\{\}}\\[0.15em]
{\footnotesize\bfseries MSC:} {\footnotesize 62F15 , 62M10 , 62H12 , 91G70 . 1cm}\\
\vfill
\end{minipage}
\newpage

% p.1283
\noindent\begin{minipage}[t][\textheight]{\textwidth}
\noindent{\bfseries\large The Bayesian nonparametric Estimation of entropy and extropy based on right censored data}\\[0.35em]
{\itshape Marani \& Saberi}\\[0.2em]
{\footnotesize First author: Marani}\\[0.45em]
\noindent\rule{\textwidth}{0.3pt}\\[0.4em]
{\footnotesize  Entropy and extropy are central measures in information theory. In this paper, Bayesian non-parametric estimators for entropy and extropy with possibly right censored data are proposed. The approach uses the beta-Stacy process as prior distribution. Examples are presented to illustrate the performance of the estimators.}\\[0.4em]
{\footnotesize\bfseries Keywords:} {\footnotesize Beta-Stacy process, entropy, extropy, non-parametric Bayesian statistics, right censored data. \textbackslash\{\}\textbackslash\{\}}\\[0.15em]
{\footnotesize\bfseries MSC:} {\footnotesize 94A17 , 62F15 , 94A15 . 1cm}\\
\vfill
\end{minipage}
\newpage

% p.1053
\noindent\begin{minipage}[t][\textheight]{\textwidth}
\noindent{\bfseries\large A Monte Carlo EM Algorithm for Fitting ARMA Models to Censored Time Series}\\[0.35em]
{\itshape Heydar Ali Mardani-Fard :h}\\[0.2em]
{\footnotesize First author: Heydar Ali Mardani-Fard :h}\\[0.45em]
\noindent\rule{\textwidth}{0.3pt}\\[0.4em]
{\footnotesize  Censored time series arise commonly in environmental and meteorological applications when measurements fall below detection limits. Fitting Gaussian ARMA models is challenging due to an irregular observed-data likelihood. park2007 proposed a simulation-based EM-like procedure that imputes censored values from their conditional truncated Gaussian distribution. While effective, this estimator often exhibits high Monte Carlo variability and unstable convergence under severe censoring. We introduce a Monte Carlo averaged variant: at each iteration, multiple independent stochastic updates are averaged, reducing estimator variance. Our method collapses to the original procedure when the number of averaged updates equals one. Simulation studies across several ARMA models show improved stability and accuracy, especially under heavy censoring.}\\[0.4em]
{\footnotesize\bfseries Keywords:} {\footnotesize Censored time series, ARMA models, Truncated multivariate normal, EM-like estimation, and Monte Carlo averaging. \textbackslash\{\}\textbackslash\{\}}\\[0.15em]
{\footnotesize\bfseries MSC:} {\footnotesize 62M10 , 83C75 , 65C05 . 1cm}\\
\vfill
\end{minipage}
\newpage

% p.1301
\noindent\begin{minipage}[t][\textheight]{\textwidth}
\noindent{\bfseries\large Expectation-Maximization Algorithm for Efficient Optimization of Multi-Response Bridge Regression}\\[0.35em]
{\itshape Rozhina Masoudi, Amir Hossein Ghatari}\\[0.2em]
{\footnotesize First author: Rozhina Masoudi}\\[0.45em]
\noindent\rule{\textwidth}{0.3pt}\\[0.4em]
{\footnotesize  Optimizing multi-response bridge regression using conventional methods like Newton-Raphson is computationally expensive and sensitive to initialization. This paper proposes a novel Expectation-Maximization (EM) algorithm for optimizing the multi-response bridge objective. By interpreting the Generalized Gaussian prior as a Gaussian Scale Mixture, we derive an iterative reweighted least squares procedure with a straightforward M-step. Numerical experiments show that EM outperforms Newton-Raphson in most cases. The two methods give similar sparsity patterns. However, EM achieves lower prediction error on average.}\\[0.4em]
{\footnotesize\bfseries Keywords:} {\footnotesize EM algorithm, Multi-response, Penalized regression, Feature selection. \textbackslash\{\}\textbackslash\{\}}\\[0.15em]
{\footnotesize\bfseries MSC:} {\footnotesize 62J05 , 62J12 \textbackslash\{\}\textbackslash\{\}}\\
\vfill
\end{minipage}
\newpage

% p.1314
\noindent\begin{minipage}[t][\textheight]{\textwidth}
\noindent{\bfseries\large E-Bayesian Estimation for Exponential Distribution under Progressive Type II Censoring with DeGroot Loss Function}\\[0.35em]
{\itshape S.M.T.K. MirMostafaee}\\[0.2em]
{\footnotesize First author: S.M.T.K. MirMostafaee}\\[0.45em]
\noindent\rule{\textwidth}{0.3pt}\\[0.4em]
{\footnotesize  This paper investigates E-Bayesian estimation for the parameter of the exponential distribution under the DeGroot loss function based on progressively Type-II censored data. We consider two scenarios: (i) when the sample size is fixed, and (ii) when the sample size is random. A simulation study verifies that a random sample size may lead to improvement in terms of estimated risk.}\\[0.4em]
{\footnotesize\bfseries Keywords:} {\footnotesize DeGroot loss function, E-Bayesian estimation, Random sample size, Progressive Censoring, and Estimated risk. \textbackslash\{\}\textbackslash\{\}}\\[0.15em]
{\footnotesize\bfseries MSC:} {\footnotesize 62F10 , 62N02 . 1cm}\\
\vfill
\end{minipage}
\newpage

% p.1329
\noindent\begin{minipage}[t][\textheight]{\textwidth}
\noindent{\bfseries\large E-Bayesian Estimation for the Lomax Distribution Under the Precautionary Loss Function Based on Record Values}\\[0.35em]
{\itshape S.M.T.K. MirMostafaee}\\[0.2em]
{\footnotesize First author: S.M.T.K. MirMostafaee}\\[0.45em]
\noindent\rule{\textwidth}{0.3pt}\\[0.4em]
{\footnotesize  This paper addresses E-Bayesian estimation of the Lomax distribution parameter within the framework of the precautionary loss function, using upper record values. The E-posterior risks of the proposed estimators are derived. A simulation study is carried out to compare and evaluate the estimators. A simulated data example is also given. The paper concludes with a discussion.}\\[0.4em]
{\footnotesize\bfseries Keywords:} {\footnotesize Precautionary loss function, E-Bayesian estimation, Upper record values, and E-posterior risk. \textbackslash\{\}\textbackslash\{\}}\\[0.15em]
{\footnotesize\bfseries MSC:} {\footnotesize 62F10 . 1cm}\\
\vfill
\end{minipage}
\newpage

% p.1253
\noindent\begin{minipage}[t][\textheight]{\textwidth}
\noindent{\bfseries\large Weighted Total Variation Distance for Copula-Based Dependency Analysis}\\[0.35em]
{\itshape Seyed Morteza Mohammadi}\\[0.2em]
{\footnotesize First author: Seyed Morteza Mohammadi}\\[0.45em]
\noindent\rule{\textwidth}{0.3pt}\\[0.4em]
{\footnotesize  Total Variation Distance (TVD) has been recently introduced as a copula-based dependency measure. Despite its useful properties, the standard TVD treats all regions of the copula space equally, which limits its ability to focus on specific dependence structures such as tail dependence. This paper proposes a weighted version of TVD, called Weighted Total Variation Distance (wTVD), which incorporates a non-negative weighting function that integrates to one. Theoretical properties, including boundedness, symmetry, and invariance under monotonic transformations, are established. A nonparametric estimator based on the Local Likelihood Probit Transformation method is presented. Simulation studies demonstrate that wTVD can effectively distinguish between different types of tail dependencies where the standard TVD fails to do so.}\\[0.4em]
{\footnotesize\bfseries Keywords:} {\footnotesize Copula, Total Variation Distance, Weighted Total Variation Distance, Tail dependence}\\[0.15em]
{\footnotesize\bfseries MSC:} {\footnotesize 62H05, 62H20. 1cm}\\
\vfill
\end{minipage}
\newpage

% p.1038
\noindent\begin{minipage}[t][\textheight]{\textwidth}
\noindent{\bfseries\large Spatial Multi-Objective Optimization for Efficient Rain Gauge Network Design in Khuzestan, Iran}\\[0.35em]
{\itshape Mohsen Mohammadzadeh mohsen}\\[0.2em]
{\footnotesize First author: Mohsen Mohammadzadeh mohsen}\\[0.45em]
\noindent\rule{\textwidth}{0.3pt}\\[0.4em]
{\footnotesize  Reliable rainfall monitoring is fundamentally linked to the strategic placement of rain gauge stations; however, many existing networks lack a systematic spatial design. This study introduces a spatial bi-objective optimization framework to evaluate the current rain gauge network in Khuzestan Province, Iran, and to define optimal configurations for future network expansion or restructuring. A novel Hybrid AMOSA\textbackslash\{\}\_NSGA-II algorithm is utilized to navigate the trade-off between two competing objectives: minimizing the root mean total error of rainfall predictions and minimizing the number of stations to reduce operational costs. The results demonstrate that optimized configurations with fewer stations consistently outperform the current network in predictive accuracy. These findings highlight the critical role of spatial optimization in developing cost-effective, high-performance rainfall monitoring systems.}\\[0.4em]
{\footnotesize\bfseries Keywords:} {\footnotesize Rain gauge station network, Multi-objective optimization, HAN algorithm, Spatial data, Kriging. \textbackslash\{\}\textbackslash\{\}}\\[0.15em]
{\footnotesize\bfseries MSC:} {\footnotesize 62H11 , 90C29 , 62P12 . 1cm}\\
\vfill
\end{minipage}
\newpage

% p.1346
\noindent\begin{minipage}[t][\textheight]{\textwidth}
\noindent{\bfseries\large An Adaptive Hybrid Active-Passive Learning Strategy for Statistical Representativeness: A Case Study of Breast Cancer Diagnosis}\\[0.35em]
{\itshape Mohammad Moradi moradi}\\[0.2em]
{\footnotesize First author: Mohammad Moradi moradi}\\[0.45em]
\noindent\rule{\textwidth}{0.3pt}\\[0.4em]
{\footnotesize Active learning reduces the cost of supervised learning by selecting informative observations for labeling, yet the resulting samples often lack statistical representativeness, leading to biased estimates of population parameters such as disease prevalence. In this study, we propose an adaptive hybrid strategy that initially employs margin-based sampling to rapidly improve predictive accuracy and then switches to simple random sampling once a predefined accuracy threshold is reached, thereby preserving the representativeness of the labeled sample. The proposed method is evaluated on the Wisconsin Breast Cancer dataset using a Random Forest classifier. Monte Carlo experimental results demonstrate that the hybrid strategy achieves predictive performance comparable to that of pure active learning while substantially improving the statistical representativeness of the labeled sample, thus providing a reliable foundation for valid statistical inference alongside high classification accuracy.\}\textbackslash\{\}\textbackslash\{\}}\\[0.4em]
{\footnotesize\bfseries Keywords:} {\footnotesize Active learning, Hybrid learning, Random Forest, Statistical inference, Breast cancer diagnosis. \textbackslash\{\}\textbackslash\{\}}\\[0.15em]
{\footnotesize\bfseries MSC:} {\footnotesize 62H30 , 62D05 , 68T05 . 1cm}\\
\vfill
\end{minipage}
\newpage

% p.1247
\noindent\begin{minipage}[t][\textheight]{\textwidth}
\noindent{\bfseries\large A Hybrid Control Chart Based on CUSUM and Double Moving Average}\\[0.35em]
{\itshape Nader Moradi-Tavalleai n}\\[0.2em]
{\footnotesize First author: Nader Moradi-Tavalleai n}\\[0.45em]
\noindent\rule{\textwidth}{0.3pt}\\[0.4em]
{\footnotesize  Quality assurance in industries is not possible without continuous monitoring and rapid detection of unwanted shifts in the process mean. In this paper, a new hybrid control chart, as DMA-CUSUM, is designed which is a combination of the Double Moving Average and CUSUM control charts. The performance of the proposed control chart is evaluated against the Double Moving Average and Moving Average-Cusum control charts using two criteria: the average run length and the median run length. The findings indicate that the proposed control chart performs superiorly in detecting small-to-moderate shifts in the process mean, particularly by choosing the appropriate width.}\\[0.4em]
{\footnotesize\bfseries Keywords:} {\footnotesize Statistical quality control, Double moving average, CUSUM chart, Average run length, Median run length, Monte Carlo simulations}\\[0.15em]
{\footnotesize\bfseries MSC:} {\footnotesize 62P30 .}\\
\vfill
\end{minipage}
\newpage

% p.1015
\noindent\begin{minipage}[t][\textheight]{\textwidth}
\noindent{\bfseries\large Covariance Estimation Strategies in Massive Spatial and Spatial-Temporal Data}\\[0.35em]
{\itshape Ali M. Mosammam}\\[0.2em]
{\footnotesize First author: Ali M. Mosammam}\\[0.45em]
\noindent\rule{\textwidth}{0.3pt}\\[0.4em]
{\footnotesize  Covariance plays very important roles in modelling and in prediction of spatial data. For instance, many statistical inferences such as maximum likelihood estimation and best linear unbiased prediction involve covariance matrices which must be inverted and this evaluation is slow when the number of observations is large, and the complexity increases even further in a space-time context. First, we addresses a Bayesian grouping-based approach, extending Vecchia's approximation and relying on the conditional independence assumption of ordered data which significantly reduces computational burden in nonparametric estimation of a nonstationary spatial covariance. Next, we propose half spectral composite likelihood for the estimation of spatial-temporal covariances which involves a spectral approach in time and a covariance function in space. The approach requires no matrix inversions and the estimators are shown to be consistent and asymptotically normal under increasing domain asymptotic. Examples are given using simulated and real data.}\\[0.4em]
{\footnotesize\bfseries Keywords:} {\footnotesize Half-spectral model, Composite Likelihood, Whittle likelihood, Vecchia approximation \textbackslash\{\}\textbackslash\{\}}\\[0.15em]
{\footnotesize\bfseries MSC:} {\footnotesize 60M30 , 62M11 , 62M30 , 62H11 . 1cm}\\
\vfill
\end{minipage}
\newpage

% p.1276
\noindent\begin{minipage}[t][\textheight]{\textwidth}
\noindent{\bfseries\large A Similarity-based Method for Saving Runs in Factorial Experiments}\\[0.35em]
{\itshape Mohsen Motavaze}\\[0.2em]
{\footnotesize First author: Mohsen Motavaze}\\[0.45em]
\noindent\rule{\textwidth}{0.3pt}\\[0.4em]
{\footnotesize  This paper presents a method for saving experimental runs during the planning stage of factorial experiments when resources for conducting a complete experimental design are limited. In this context, the method investigates the number of coincidences between experimental runs as a measure of similarity. The real advancement of this approach lies in its emphasis on the dynamic interactions between the runs, rather than just focusing on the relationships among the factors represented as columns. Despite its simplicity, the method shows impressive performance in numerical evaluations.}\\[0.4em]
{\footnotesize\bfseries Keywords:} {\footnotesize coincidence measure, experimental run, Hamming distance, negligible contrasts, similarity. \textbackslash\{\}\textbackslash\{\}}\\[0.15em]
{\footnotesize\bfseries MSC:} {\footnotesize 62K15 . 1cm}\\
\vfill
\end{minipage}
\newpage

% p.1214
\noindent\begin{minipage}[t][\textheight]{\textwidth}
\noindent{\bfseries\large BehaveSeg: A Behavior-Aware Multi-Algorithm Framework for Intelligent Customer Segmentation}\\[0.35em]
{\itshape Mobina Mozayan}\\[0.2em]
{\footnotesize First author: Mobina Mozayan}\\[0.45em]
\noindent\rule{\textwidth}{0.3pt}\\[0.4em]
{\footnotesize  Customer segmentation remains a cornerstone of personalized marketing, yet conventional methods frequently fall short in capturing dynamic behavioral patterns due to their reliance on static demographic or transactional data. This study introduces BehaveSeg , a behavior-aware pipeline that integrates rich behavioral feature engineering, Principal Component Analysis (PCA) dimensionality reduction, and a comprehensive multi-algorithm clustering approach. Seven widely used clustering algorithms were systematically evaluated using Silhouette Score, Davies--Bouldin Index, and Calinski--Harabasz Index, with the first seven PCA components preserving \textbackslash\{\}(92.2\textbackslash\{\}}\\[0.4em]
{\footnotesize\bfseries Keywords:} {\footnotesize Customer Segmentation, Clustering Analysis, Behavioral Analytics, Unsupervised Learning \textbackslash\{\}\textbackslash\{\}}\\[0.15em]
{\footnotesize\bfseries MSC:} {\footnotesize 62H30 , 68T05 , 91B42 , 62-07 . 1cm}\\
\vfill
\end{minipage}
\newpage

% p.1205
\noindent\begin{minipage}[t][\textheight]{\textwidth}
\noindent{\bfseries\large Construction and Analysis of a Nonstandard Finite Difference Scheme for Stochastic Partial Differential Equations}\\[0.35em]
{\itshape Mehran Namjoo}\\[0.2em]
{\footnotesize First author: Mehran Namjoo}\\[0.45em]
\noindent\rule{\textwidth}{0.3pt}\\[0.4em]
{\footnotesize  This paper is devoted to the development and comprehensive analysis of a nonstandard finite difference (NSFD) scheme for It\textbackslash\{\}\textasciicircum{} o stochastic partial differential equations (SPDEs). The qualitative properties of the proposed scheme, including consistency, stability, and convergence, are thoroughly investigated. Numerical simulations are presented to demonstrate the accuracy and computational efficiency of the proposed scheme.}\\[0.4em]
{\footnotesize\bfseries Keywords:} {\footnotesize Nonstandard finite difference scheme, consistency, stability, convergence, stochastic partial differential equations. \textbackslash\{\}\textbackslash\{\}}\\[0.15em]
{\footnotesize\bfseries MSC:} {\footnotesize 65M06 , 65M12 , 60H15 . 1cm}\\
\vfill
\end{minipage}
\newpage

% p.1262
\noindent\begin{minipage}[t][\textheight]{\textwidth}
\noindent{\bfseries\large Nonparametric Bayesian Optimal Design for Exponential-Shift Regression Model with stick breaking method and uniform distribution as the base measure}\\[0.35em]
{\itshape Anita Abdollahi Nanvapisheh, Atefeh Abdollahi Nanvapisheh}\\[0.2em]
{\footnotesize First author: Anita Abdollahi Nanvapisheh}\\[0.45em]
\noindent\rule{\textwidth}{0.3pt}\\[0.4em]
{\footnotesize  red In this study, we initially present the shifted exponential regression model and then derive its corresponding optimal experimental designs by employing the Dirichlet Process (DP) as a prior distribution over the parameter space. The resulting optimal design problem is formulated as a nonlinear optimization task, and the final solutions are obtained using the Rsolnp package implemented in the R statistical computing environment.}\\[0.4em]
{\footnotesize\bfseries Keywords:} {\footnotesize Exponential-Shift Regression, Lifetime Data Analysis, Reliability Modeling, Nonparametric Bayesian Optimal Design, Dirichlet Process. \textbackslash\{\}\textbackslash\{\}}\\[0.15em]
{\footnotesize\bfseries MSC:} {\footnotesize 60E05 , 60Exx . 1cm}\\
\vfill
\end{minipage}
\newpage

% p.1068
\noindent\begin{minipage}[t][\textheight]{\textwidth}
\noindent{\bfseries\large Bayesian Inference for Matrix-Variate Gaussian Distributions: Theory, Computation, and Applications}\\[0.35em]
{\itshape Fariba Nasirzadeh}\\[0.2em]
{\footnotesize First author: Fariba Nasirzadeh}\\[0.45em]
\noindent\rule{\textwidth}{0.3pt}\\[0.4em]
{\footnotesize  We develop a fully Bayesian framework for matrix-variate Gaussian distributions with Kronecker-structured covariance , including conjugate priors, Gibbs sampling full conditionals, and sparsity-inducing extensions for high dimensions. The inherent scaling identifiability issue ( ) = (c ) (c\textasciicircum{} -1 ) is addressed via the constraint tr ( )=p , though simulations reveal asymmetry when p q . Increasing sample size from n=30 to n=100 reduces relative Frobenius error by 55\textbackslash\{\}}\\[0.4em]
{\footnotesize\bfseries Keywords:} {\footnotesize Bayesian inference, Gibbs sampling, Kronecker covariance, Matrix-variate Gaussian}\\[0.15em]
{\footnotesize\bfseries MSC:} {\footnotesize 62F15, 62H12, 62J07. 1cm}\\
\vfill
\end{minipage}
\newpage

% p.1035
\noindent\begin{minipage}[t][\textheight]{\textwidth}
\noindent{\bfseries\large On the Zero-Inflated Geometric INAR(1 Process with Random Coefficient}\\[0.35em]
{\itshape R. Nasirzadeh}\\[0.2em]
{\footnotesize First author: R. Nasirzadeh}\\[0.45em]
\noindent\rule{\textwidth}{0.3pt}\\[0.4em]
{\footnotesize  Count time series with overdispersion, excess zeros, and autocorrelation frequently arise in hydrology, epidemiology, and climatology. However, INAR processes with random coefficients have received limited attention, and non-random coefficient models often fail due to ignored random effects. Moreover, neglecting zero inflation causes overdispersion and biased estimates. This paper examines the zero-inflated geometric INAR(1) process with random coefficient ( ZIGINAR \_ RC (1) ), synthesizing its probabilistic properties, Whittle estimation, sieve bootstrap methods, and Bayesian forecasting. Simulation results show Bayesian prediction achieves superior accuracy, while sieve bootstrap offers good performance with reasonable efficiency. Practical applicability is demonstrated using animal health submission data.}\\[0.4em]
{\footnotesize\bfseries Keywords:} {\footnotesize Count time series, Zero-inflation, Random coefficient, Estimation methods, Forecasting}\\[0.15em]
{\footnotesize\bfseries MSC:} {\footnotesize 62M10, 62F10, 62M20. 1cm}\\
\vfill
\end{minipage}
\newpage

% p.1311
\noindent\begin{minipage}[t][\textheight]{\textwidth}
\noindent{\bfseries\large Bayesian Structure Learning of Gene Expression Networks via Laplace-Approximated Scoring: An Application to TCGA Breast Cancer Data}\\[0.35em]
{\itshape Nazari et al}\\[0.2em]
{\footnotesize First author: Nazari}\\[0.45em]
\noindent\rule{\textwidth}{0.3pt}\\[0.4em]
{\footnotesize  Breast cancer is a heterogeneous disease driven by complex molecular interactions. Identifying direct dependencies between continuous gene expression profiles and discrete clinical outcomes is critical to robust biomarker discovery. In this study, we present an empirical application of a recently developed Bayesian directed acyclic graph (DAG) learning framework to infer probabilistic dependency networks from the TCGA Breast Cancer dataset. The proposed framework utilizes a DAG-Probit model, treating the binary clinical status as a thresholded latent Gaussian variable. By leveraging Laplace approximations to efficiently estimate node-wise marginal likelihoods and employing Markov chain Monte Carlo (MCMC) sampling for high-dimensional structure exploration, the model successfully infers a sparse, clinically relevant regulatory network. The application identifies key driver genes---such as COL11A1 , PCOLCE2 , and SFRP1 ---with a posterior edge inclusion probability ( 1.00 ), demonstrating the effectiveness of the framework for biological interpretation and mechanistic oncology studies.}\\[0.4em]
{\footnotesize\bfseries Keywords:} {\footnotesize Bayesian Structure Learning, Breast Cancer Data, DAG-Probit Model, Gene Regulatory Networks, Mixed Data. \textbackslash\{\}\textbackslash\{\}}\\[0.15em]
{\footnotesize\bfseries MSC:} {\footnotesize 62P10 , 62F15 , 62H22 . 1cm}\\
\vfill
\end{minipage}
\newpage

% p.1239
\noindent\begin{minipage}[t][\textheight]{\textwidth}
\noindent{\bfseries\large Beyond 1D Heuristics: Temporal-Causal Calibration and Multivariate Non-Parametric Hypothesis Testing for Uncertainty Quantification in Retrieval-Augmented Language Models}\\[0.35em]
{\itshape Sepehr Nazeri, Dr. Ferdos Gorji}\\[0.2em]
{\footnotesize First author: Sepehr Nazeri}\\[0.45em]
\noindent\rule{\textwidth}{0.3pt}\\[0.4em]
{\footnotesize  Most methods for measuring how confident a language model should be rely on a single number, predictive entropy, which is too simple to capture what really happens when a model's built-in knowledge conflicts with facts from an external source. We propose the Temporal-Causal Calibration (TCC) framework, which measures parametric stubbornness - how strongly a model clings to old knowledge when new evidence disagrees with it, using tools from causal inference. We then use statistical hypothesis tests to check whether the model's uncertainty scores are trustworthy. Tested on two small, locally run models, TCC-UQ achieves an AUROC of 0.4895 and an Expected Calibration Error (ECE) of 0.1337 , using only 2.1\textbackslash\{\},MB of extra memory. Kolmogorov-Smirnov tests ( D = 0.713 , p 4.90 10\textasciicircum{} -37 ) and paired t -tests ( t = -64.48 , p 9.66 10\textasciicircum{} -111 ) confirm that TCC-UQ produces a statistically distinct confidence distribution from the uncalibrated baseline. The results demonstrate that causal grounding of stubbornness measurement yields interpretable, statistically verifiable uncertainty signals, while also revealing important limitations that motivate future work.\textbackslash\{\}\textbackslash\{\}}\\[0.4em]
{\footnotesize\bfseries Keywords:} {\footnotesize Retrieval-Augmented Generation; Uncertainty Quantification(UQ); Calibration; Non-Parametric Testing; Causal Inference.\textbackslash\{\}\textbackslash\{\}}\\[0.15em]
\vfill
\end{minipage}
\newpage

% p.1354
\noindent\begin{minipage}[t][\textheight]{\textwidth}
\noindent{\bfseries\large High-Resolution Spectral Estimation of Univariate and Multivariate Stationary Time Series}\\[0.35em]
{\itshape A. R. Nematollahi}\\[0.2em]
{\footnotesize First author: A. R. Nematollahi}\\[0.45em]
\noindent\rule{\textwidth}{0.3pt}\\[0.4em]
{\footnotesize  This study evaluates spectral density estimation techniques for stationary time series, with an emphasis on high-resolution methods. Classical approaches often have difficulty resolving closely-spaced frequency components when data is limited. Through controlled simulations, we compare these traditional techniques against high-resolution estimators and show that the latter consistently deliver clearer, more distinct frequency estimates. The methods are also extended to multivariate time series, demonstrating their practical value in complex signal characterization.}\\[0.4em]
{\footnotesize\bfseries Keywords:} {\footnotesize Stationary time series, Spectral density, Capon's estimate, High resolution estimate, Block-Toeplitz matrix, Periodogram. \textbackslash\{\}\textbackslash\{\}}\\[0.15em]
{\footnotesize\bfseries MSC:} {\footnotesize 62M15 , 62G07 , 15B05 . 1cm}\\
\vfill
\end{minipage}
\newpage

% p.1072
\noindent\begin{minipage}[t][\textheight]{\textwidth}
\noindent{\bfseries\large Intelligent Anomaly Detection with Deep Reinforcement Learning for Industrial Inspection}\\[0.35em]
{\itshape Amirhossein Khadivi Noghredeh, Abdollah Safari}\\[0.2em]
{\footnotesize First author: Amirhossein Khadivi Noghredeh}\\[0.45em]
\noindent\rule{\textwidth}{0.3pt}\\[0.4em]
{\footnotesize  This work tackles pixel-level anomaly detection in industrial inspection with scarce defective data. We introduce a semi-supervised deep reinforcement learning framework enhancing detection via an adaptive neural batch sampler and robust autoencoder, balancing exploration and exploitation for subtle anomalies. MVTec AD evaluations show superior performance (AUC: 0.949 0.03, F1\_ : 0.561 0.20), showing favorable results compared to structurally comparable lightweight methods, and standard deviations confirm robustness despite the increased mean accuracy. This offers a scalable, statistically robust solution for manufacturing quality assurance.}\\[0.4em]
{\footnotesize\bfseries Keywords:} {\footnotesize Anomaly Detection, Reinforcement Learning, Outlier Detection, Computer Vision, Semi-Supervised Learning, Adaptive Batch Sampling \textbackslash\{\}\textbackslash\{\}}\\[0.15em]
{\footnotesize\bfseries MSC:} {\footnotesize 68T07 , 68T45 , 62H30 . 1cm}\\
\vfill
\end{minipage}
\newpage

% p.1209
\noindent\begin{minipage}[t][\textheight]{\textwidth}
\noindent{\bfseries\large Deep Learning--Based Competing Risks Survival Analysis for Predicting Cause-Specific Mortality in Heart Failure Patients}\\[0.35em]
{\itshape Solmaz Norouzi , Orcid: 0000-0001-5805-6072}\\[0.2em]
{\footnotesize First author: Solmaz Norouzi}\\[0.45em]
\noindent\rule{\textwidth}{0.3pt}\\[0.4em]
{\footnotesize  Heart failure (HF) patients face multiple competing causes of death, making standard survival models prone to bias. We developed a deep learning--based competing risks (DLCR) model to predict cause-specific mortality in a retrospective cohort of 364 HF patients from a tertiary cardiovascular center in Iran. After a five-year follow-up, 160 patients died of HF and 97 of other causes. The DLCR model achieved a concordance index (C-index) of 0.71 for HF-related mortality and 0.81 for non-HF mortality, with an integrated Brier score (IBS) of 0.16 and 0.09, respectively. The model captured distinct temporal risk patterns: HF-related deaths dominated the first 35 months, after which other-cause mortality gradually exceeded HF mortality. Our findings demonstrate that explicitly accounting for competing risks via deep learning improves risk stratification accuracy. The DLCR framework offers a promising tool for personalized prognosis in HF, although external validation in multi-center cohorts is needed before clinical deployment.}\\[0.4em]
{\footnotesize\bfseries Keywords:} {\footnotesize Deep learning, Competing risks, Survival analysis, Heart failure, Cause-specific mortality, Predictive modeling}\\[0.15em]
{\footnotesize\bfseries MSC:} {\footnotesize 62N01, 62N02, 92C50. 1cm}\\
\vfill
\end{minipage}
\newpage

% p.1163
\noindent\begin{minipage}[t][\textheight]{\textwidth}
\noindent{\bfseries\large A Hybrid FCNN--LSTM Framework for Sequential Prediction with Application to Air Pollution Forecasting}\\[0.35em]
{\itshape Abbas Pak}\\[0.2em]
{\footnotesize First author: Abbas Pak}\\[0.45em]
\noindent\rule{\textwidth}{0.3pt}\\[0.4em]
{\footnotesize  Accurate prediction of sequential data is challenging due to complex nonlinear relationships and temporal dependencies. Fully Connected Neural Networks (FCNNs) effectively capture nonlinear feature interactions but cannot explicitly model temporal patterns, whereas Long Short-Term Memory (LSTM) networks are suitable for sequential learning but may inadequately represent complex feature relationships. This study proposes a hybrid FCNN--LSTM framework that combines nonlinear feature extraction with temporal sequence modeling. The FCNN component learns latent feature representations, while the LSTM component captures dynamic temporal dependencies. The proposed framework is evaluated on an air pollution prediction task using real-world environmental data. Experimental results show that the hybrid model consistently outperforms standalone FCNN and LSTM models across multiple evaluation metrics, demonstrating superior predictive accuracy and generalization capability.}\\[0.4em]
{\footnotesize\bfseries Keywords:} {\footnotesize Fully Connected Neural Networks, Long Short-Term Memory Networks, Sequential Data, and Temporal Dependency \textbackslash\{\}\textbackslash\{\}}\\[0.15em]
{\footnotesize\bfseries MSC:} {\footnotesize 92B20 , 62M20 . 1cm}\\
\vfill
\end{minipage}
\newpage

% p.1254
\noindent\begin{minipage}[t][\textheight]{\textwidth}
\noindent{\bfseries\large The Critical Role of Feature Engineering in Time-Series Prediction: A Case Study on Air Pollution Forecasting}\\[0.35em]
{\itshape Pak and Kaviani}\\[0.2em]
{\footnotesize First author: Pak}\\[0.45em]
\noindent\rule{\textwidth}{0.3pt}\\[0.4em]
{\footnotesize  Feature engineering plays a crucial role in improving the predictive performance of machine learning models for time-series forecasting. This study investigates the impact of engineered temporal features on air pollution prediction using a dataset containing atmospheric pollutants and meteorological variables. A comprehensive feature engineering framework was developed, including lag variables, rolling statistics, cyclical time encodings, and pollutant--weather interaction features. The effectiveness of the proposed features was evaluated using the Random Forest and Light Gradient Boosting Machine (LightGBM) models. Forecasting performance was compared before and after feature engineering for multiple pollutants, including CO, O3, NO2, and SO2. Results showed substantial improvements across all prediction tasks, demonstrating the ability of the derived features to capture temporal dependencies, seasonal patterns, and nonlinear environmental interactions. The results highlight the pivotal role of feature engineering in time-series forecasting and demonstrate that substantial accuracy improvements can be achieved without modifying the underlying learning algorithm.}\\[0.4em]
{\footnotesize\bfseries Keywords:} {\footnotesize Time-series Forecasting, Feature engineering, Random Forest, and Light Gradient Boosting Machine. \textbackslash\{\}\textbackslash\{\}}\\[0.15em]
{\footnotesize\bfseries MSC:} {\footnotesize 92B20 , 62M20 . 1cm}\\
\vfill
\end{minipage}
\newpage

% p.1260
\noindent\begin{minipage}[t][\textheight]{\textwidth}
\noindent{\bfseries\large Extropy of the Residual Lifetime of Alive Components in Coherent Systems}\\[0.35em]
{\itshape Zohreh Pakdaman, Reza Alizadeh Noughabi}\\[0.2em]
{\footnotesize First author: Zohreh Pakdaman}\\[0.45em]
\noindent\rule{\textwidth}{0.3pt}\\[0.4em]
{\footnotesize  This research investigates the uncertainty associated with the remaining lifetimes of alive components in coherent systems after system failure. Although the system may no longer function, some components may still remain alive, and the uncertainty in their remaining lifetimes can be meaningfully quantified. To this end, we employ extropy as a measure of uncertainty for the remaining lifetimes of alive components. By combining the system signature with conditional properties of order statistics, we develop a mathematical framework for this purpose. Focusing on the heavy-tailed Lomax distribution, our results provide insight into how the system structure affects the uncertainty of the remaining lifetimes of alive components after system failure.}\\[0.4em]
{\footnotesize\bfseries Keywords:} {\footnotesize Extropy, Coherent System, Order statistics, System signature, Residual lifetimes of live components. \textbackslash\{\}\textbackslash\{\}}\\[0.15em]
{\footnotesize\bfseries MSC:} {\footnotesize 62N05 ; 62F10 . 1cm}\\
\vfill
\end{minipage}
\newpage

% p.1340
\noindent\begin{minipage}[t][\textheight]{\textwidth}
\noindent{\bfseries\large Graphical Finite-Population Sampling with Adaptive Monte Carlo Tree Search}\\[0.35em]
{\itshape Bardia Panahbehagh}\\[0.2em]
{\footnotesize First author: Bardia Panahbehagh}\\[0.45em]
\noindent\rule{\textwidth}{0.3pt}\\[0.4em]
{\footnotesize  Graphical finite-population sampling represents prescribed first-order inclusion probabilities by bars, or unions of bar segments, on the unit interval. A horizontal uniform random line induces a sample, while the geometry of overlaps determines the second- and higher-order inclusion probabilities. Starting from the fixed-size graphical construction equivalent to Madow's systematic design, exchanges of interchangeable segments generate alternative fixed-size designs without changing the prescribed first-order probabilities. This paper combines that feasible design generator with an adaptive population-based Monte Carlo tree search. The implementation uses upper-confidence tree selection, quality-adaptive mutation scales, parallel two-scale rollouts, a safety cutoff, and restarts from the best design after stagnation. A tuning experiment based on the exact Narain--Horvitz--Thompson variance indicates that the restart and rollout mechanisms materially affect the quality of the best graphical design found.}\\[0.4em]
{\footnotesize\bfseries Keywords:} {\footnotesize Finite-population sampling, graphical sampling, inclusion probabilities, Monte Carlo tree search, Narain--Horvitz--Thompson estimator. \textbackslash\{\}\textbackslash\{\}}\\[0.15em]
{\footnotesize\bfseries MSC:} {\footnotesize 62D05 , 62K05 , 90C59 . \textbackslash\{\}\textbackslash\{\}}\\
\vfill
\end{minipage}
\newpage

% p.1043
\noindent\begin{minipage}[t][\textheight]{\textwidth}
\noindent{\bfseries\large Machine Learning Methods for Bivariate Integer-Valued Time-Series Data}\\[0.35em]
{\itshape Mahnaz Pezeshkian, Maryam Sharafi}\\[0.2em]
{\footnotesize First author: Mahnaz Pezeshkian}\\[0.45em]
\noindent\rule{\textwidth}{0.3pt}\\[0.4em]
{\footnotesize  In this paper, we approach the problem through a machine learning framework, focusing on time series data particularly integer-valued observations without resorting to classical statistical analysis. First, we introduce four machine learning techniques. Then, the procedure will be examined, and assessing their predictive accuracy using mean absolute error. In application, this study analyze a count time series dataset using machine learning techniques. Our results highlight the strengths of the approach and provide insights for model selection in count time series applications.\textbackslash\{\}\textbackslash\{\}}\\[0.4em]
{\footnotesize\bfseries Keywords:} {\footnotesize Machine learning techniques, Cross-validation, Bivariate integer-valued autoregressive models, Over-dispersion. \textbackslash\{\}\textbackslash\{\}}\\[0.15em]
{\footnotesize\bfseries MSC:} {\footnotesize Machine Learning 1cm}\\
\vfill
\end{minipage}
\newpage

% p.1266
\noindent\begin{minipage}[t][\textheight]{\textwidth}
\noindent{\bfseries\large Joint Modeling of Zero-Inflated Longitudinal Count Data and Survival Outcomes with Cure Fraction}\\[0.35em]
{\itshape Sahar Souri Pilangorgi, Soheila Khodakarim}\\[0.2em]
{\footnotesize First author: Sahar Souri Pilangorgi}\\[0.45em]
\noindent\rule{\textwidth}{0.3pt}\\[0.4em]
{\footnotesize  Background: Longitudinal data with excess zeros and survival data with a cured fraction are common in medical research, often exhibiting interdependence that complicates analysis. Joint modeling of these data types can address this interdependence, potentially improving estimation accuracy over independent models. This study introduces a joint model that integrates a zero-inflated count model for the longitudinal process with a mixture cure model for the survival process and investigates the performance of joint models compared to independent models under varying conditions of zero-inflation, cure rates, and sample sizes. Methods: We analyzed 500 datasets with sample sizes of 500 and 800 as part of a simulation study that included 32 scenarios. Zero-Inflated Poisson (ZIP) and Zero-Inflated Negative Binomial (ZINB) distributions were used to generate longitudinal count data, with zero-inflation rates of 43\textbackslash\{\} Results: Joint models consistently outperformed independent models, showing lower bias across most parameters. Joint modeling of ZINB longitudinal count data and survival data with a cure fraction (JZINBC model) exhibited less bias than joint modeling of ZIP longitudinal count data and survival data with a cure fraction (JZIPC model) in high zero-inflation settings due to better handling of overdispersion, while JZIPC model performed better with lower zero-inflation (43\textbackslash\{\} Conclusion: Joint modeling of zero-inflated longitudinal count data and survival data with a cure fraction (JZIC model) provides more accurate parameter estimates than independent models, particularly in complex scenarios with high zero-inflation and cure rates. These findings highlight the importance of joint modeling in medical research involving interdependent longitudinal and survival outcomes.}\\[0.4em]
{\footnotesize\bfseries Keywords:} {\footnotesize Zero-inflated, longitudinal count data, mixture cure model, joint model. \textbackslash\{\}\textbackslash\{\}}\\[0.15em]
{\footnotesize\bfseries MSC:} {\footnotesize 99X99 , 99X99 , 99X99 . 1cm}\\
\vfill
\end{minipage}
\newpage

% p.1045
\noindent\begin{minipage}[t][\textheight]{\textwidth}
\noindent{\bfseries\large Optimal lot reinspection based on truncated life test for logistic Rayleigh distribution}\\[0.35em]
{\itshape Mehran Naghizadeh Qomi}\\[0.2em]
{\footnotesize First author: Mehran Naghizadeh Qomi}\\[0.45em]
\noindent\rule{\textwidth}{0.3pt}\\[0.4em]
{\footnotesize  Single sampling scheme is a conventional method for evaluating the equality of lots of products. Repetitive sampling plans are introduced in order to reduce the size of samples drawn from the lot. The product lifetime follows the logistic Rayleigh distribution. Best repetitive sampling plans are determined by minimizing the sampling effort using integer nonlinear programming. In terms of average sample number, optimal repetitive sampling plans outperform single sampling plans. A practical application is included for illustrative purposes.\textbackslash\{\}\textbackslash\{\}}\\[0.4em]
{\footnotesize\bfseries Keywords:} {\footnotesize Quality control, Repetitive sampling plans, Average sample number, Producer and Consumer risks. \textbackslash\{\}\textbackslash\{\}}\\[0.15em]
{\footnotesize\bfseries MSC:} {\footnotesize 62N05 , 62P30 . 1cm}\\
\vfill
\end{minipage}
\newpage

% p.1092
\noindent\begin{minipage}[t][\textheight]{\textwidth}
\noindent{\bfseries\large Moving Extreme Ranked Set Sampling for Imbalanced Binary Classification: A Logistic Regression Perspective}\\[0.35em]
{\itshape R.Yousefi, H.Doosti,H.Esmaily}\\[0.2em]
{\footnotesize First author: R.Yousefi}\\[0.45em]
\noindent\rule{\textwidth}{0.3pt}\\[0.4em]
{\footnotesize  Logistic regression (LR) is widely used for binary classification, yet its performance is often compromised by class imbalance and inefficient sampling. Moving extreme ranked set sampling (MERSS) offers a potential solution by utilizing auxiliary information to target extreme observations where rare events reside, but its impact on LR classification has not been systematically evaluated. This study compares SRS and MERSS for LR classification under imbalanced conditions (\textbackslash\{\}(r = 0.2, 0.3, 0.4\textbackslash\{\})), varying sample sizes (\textbackslash\{\}(n = 120, 210, 300\textbackslash\{\})), correlations (\textbackslash\{\}( = 0.3, 0.5, 0.7\textbackslash\{\})), and set sizes (\textbackslash\{\}(k = 3, 5, 10\textbackslash\{\})). Results demonstrate that MERSS consistently outperforms SRS, particularly for minority class identification. The largest gains were observed under the most severe imbalance (\textbackslash\{\}(r = 0.2\textbackslash\{\})), where MERSS improved Precision by up to 15\textbackslash\{\}}\\[0.4em]
{\footnotesize\bfseries Keywords:} {\footnotesize Logistic regression, Moving extreme ranked set sampling (MERSS), Imbalanced classification \textbackslash\{\}\textbackslash\{\}}\\[0.15em]
{\footnotesize\bfseries MSC:} {\footnotesize 62D05 , 62J12 , 62P10 . 1cm}\\
\vfill
\end{minipage}
\newpage

% p.1350
\noindent\begin{minipage}[t][\textheight]{\textwidth}
\noindent{\bfseries\large S. Nezamdoust}\\[0.35em]
{\itshape Soghra Rahimi; and Farzad Eskandari}\\[0.2em]
{\footnotesize First author: Soghra Rahimi}\\[0.45em]
\noindent\rule{\textwidth}{0.3pt}\\[0.4em]
{\footnotesize This paper proposes a novel enhancement to semi-parametric models by integrating shrinkage to improve the predictive accuracy and performance of Ridge and ecpc models. By employing the EM algorithm, updated parameter estimates are derived, demonstrating significant error reductions in metrics such as Mean Squared Error (MSE , Sum of Squared Errors (SSE , and Geometric Mean Squared Error (GMSE . Extensive simulation studies reveal consistent improvements in model performance, with shrinkage enhancing accuracy and robustness across various scenarios. The findings highlight the potential of shrinkage in refining semi-parametric models, offering a more accurate framework for statistical predictions. This study paves the way for future research into applying shrinkage techniques to diverse data types and model structures, setting a benchmark for model optimization in statistical analysis}\\[0.4em]
\vfill
\end{minipage}
\newpage

% p.1335
\noindent\begin{minipage}[t][\textheight]{\textwidth}
\noindent{\bfseries\large Efficient MCP-Penalized Variable Selection for Time-Dependent Cox Models with Counting-Process Data}\\[0.35em]
{\itshape Rohollah Ramezani, Mohammad Reza Rabiei rabiei}\\[0.2em]
{\footnotesize First author: Rohollah Ramezani}\\[0.45em]
\noindent\rule{\textwidth}{0.3pt}\\[0.4em]
{\footnotesize  This paper proposes an MCP-penalized approach for variable selection in time-dependent Cox regression models with counting-process data. The proposed algorithm enables efficient estimation in moderate- and high-dimensional settings. Simulation results show that MCP substantially reduces false negatives and is computationally more efficient than LASSO. An analysis of the Mayo Clinic PBC dataset further demonstrates that MCP achieves predictive performance comparable to LASSO while requiring less computation time and producing less shrinkage of important covariate effects.}\\[0.4em]
{\footnotesize\bfseries Keywords:} {\footnotesize Time-Dependent Cox Model, Minimax Concave Penalty (MCP), Variable Selection, and Counting-Process Data. \textbackslash\{\}\textbackslash\{\}}\\[0.15em]
{\footnotesize\bfseries MSC:} {\footnotesize 62N02 , 62J07 , 62J12 . 1cm}\\
\vfill
\end{minipage}
\newpage

% p.1161
\noindent\begin{minipage}[t][\textheight]{\textwidth}
\noindent{\bfseries\large Bayesian Local Sensitivity to Non-ignorability with General Missingness of Outcomes and Covariates}\\[0.35em]
{\itshape Elahe Momeni Roochi e}\\[0.2em]
{\footnotesize First author: Elahe Momeni Roochi e}\\[0.45em]
\noindent\rule{\textwidth}{0.3pt}\\[0.4em]
{\footnotesize  A major concern in statistical analysis is missing data. Analysis frequently operates under the presumption that missingness is ignorable when it occurs in the study. It is critical to evaluate the effect of missingness on the most important Bayesian inferences because ignorability is an untestable assumption in the absence of data augmentation. Previously, in a Bayesian framework, the influence of missing outcomes' non-ignorability was studied using local sensitivity indices. In this paper, for the first time, the Bayesian index of local sensitivity to non-ignorability for data exposed to general missingness in outcome and covariates will be extended. My results indicate that the Bayesian estimates of parameters of interest via local sensitivity approximation are more accurate than complete-case analysis and ignorable model. I will evaluate the features of my proposed indices through several simulated studies.}\\[0.4em]
{\footnotesize\bfseries Keywords:} {\footnotesize Local Sensitivity to Non-ignorability, General Missingness, Bayesian Approach, and Local Sensitivity Indices.\textbackslash\{\}\textbackslash\{\}}\\[0.15em]
{\footnotesize\bfseries MSC:} {\footnotesize 99X99 , 99X99 , 99X99 . 1cm}\\
\vfill
\end{minipage}
\newpage

% p.1330
\noindent\begin{minipage}[t][\textheight]{\textwidth}
\noindent{\bfseries\large On the Asymptotic Properties of Liu Estimator}\\[0.35em]
{\itshape M. Roozbeh}\\[0.2em]
{\footnotesize First author: M. Roozbeh}\\[0.45em]
\noindent\rule{\textwidth}{0.3pt}\\[0.4em]
{\footnotesize  In traditional regression analysis, ordinary least-squares estimation is the optimal approach for deriving regression weights, provided that the fundamental assumptions are satisfied. Nevertheless, if the data fail to meet certain assumptions, the findings may be deceptive. Outliers particularly contravene the assumption of regularly distributed residuals in least-squares regression. Therefore, it is essential to employ estimating techniques specifically formulated to address these issues. In this paper, the Liu estimator and its robust type are introduced in the semiparametric linear model when the error terms are not independent and some linear constraints are assumed to combat the presence of multicollinearity and outlier challenges. In the literature on Liu methodology, the obtaining of biased parameter plays an important role in model prediction and estimation. Based on a set of observations with a sample size of n , and the integer trimming parameter h n , the LTS estimator computes the hyperplane that minimizes the sum of the least h squared residuals. Even though LTS estimator is statistically more effective than the widely used least median squares (LMS) estimate, it is less complicated computationally than LMS. Hereby we propose the robust estimators for the biasing parameter of Liu estimation based on the improved LTS obtained by semidefinite programming approach. Asymptotic normality of proposed robust Liu estimators under some conditions is also proved and a robust test is given for testing the symmetry hypothesis H\_0: R = r . It is shown that these estimators perform better than classical Liu estimator. For our proposal, via a Monte-Carlo simulation and a real data example, performance of the Liu type of robust estimators are compared with the classical ones in restricted partially linear models.}\\[0.4em]
{\footnotesize\bfseries Keywords:} {\footnotesize Breakdown point, Least trimmed squares estimator, Linear restriction, Multicollinearity, Outlier, Semidefinite programming, Semiparametric regression models \textbackslash\{\}\textbackslash\{\}}\\[0.15em]
{\footnotesize\bfseries MSC:} {\footnotesize 62J20 , , 62G35 . 1cm}\\
\vfill
\end{minipage}
\newpage

% p.1210
\noindent\begin{minipage}[t][\textheight]{\textwidth}
\noindent{\bfseries\large Multivariate Tail Conditional Expectation Under a Linear Constraint: A New Risk Measure for Aggregated Exposures}\\[0.35em]
{\itshape Roohollah Roozegar}\\[0.2em]
{\footnotesize First author: Roohollah Roozegar}\\[0.45em]
\noindent\rule{\textwidth}{0.3pt}\\[0.4em]
{\footnotesize  Most multivariate tail risk measures, such as the multivariate tail conditional expectation (MTCE), assume that all risk factors exceeds their marginal quantiles simultaneously. But, in real word, this assumption sometimes is unrealistic, since systemic events are typically generated by a linear combination of risk factors crossing a critical threshold. So, in this paper, we propose a new risk measure called the multivariate tail conditional expectation under a linear constraint (MTCE-LC). This new risk measure captures the expected risk conditional on a weighted sum of risk factors exceeding a high quantile, so that weakness in one risk factor can be offset by strength in another. This measure can be a good tool for stress testing and capital allocation in systemic risk management. We present an exact expression for the proposed measure under multivariate normal risk factors and also, we carry out a Monte Carlo simulation study to show the accuracy of the theoretical results.}\\[0.4em]
{\footnotesize\bfseries Keywords:} {\footnotesize Tail conditional expectation, Multivariate risk measure, Linear constraint, Systemic risk, Multivariate normal distribution. \textbackslash\{\}\textbackslash\{\}}\\[0.15em]
{\footnotesize\bfseries MSC:} {\footnotesize 62H10 , 60E05 . 1cm}\\
\vfill
\end{minipage}
\newpage

% p.1235
\noindent\begin{minipage}[t][\textheight]{\textwidth}
\noindent{\bfseries\large A Fuzzy Spatio-Temporal Logistic Regression Model with Latent Skew- Random Field}\\[0.35em]
{\itshape Mohammad Mehdi Saber}\\[0.2em]
{\footnotesize First author: Mohammad Mehdi Saber}\\[0.45em]
\noindent\rule{\textwidth}{0.3pt}\\[0.4em]
{\footnotesize  Conventional spatio-temporal logistic regression models relying on Gaussian latent fields often impose restrictive assumptions of symmetry and light tails. Furthermore, they struggle to incorporate epistemic uncertainty, such as measurement imprecision, without relying on unverifiable distributional assumptions. To address these limitations, we propose a novel fuzzy spatio-temporal logistic regression model featuring a latent skew- t random field. Our approach leverages a Gamma scale-mixture to capture aleatory asymmetry and heavy tails, while employing fuzzy numbers to non-parametrically represent imprecision in covariates and parameters. We develop a Fuzzy Monte Carlo Expectation-Maximization algorithm for parameter estimation and establish its consistency. Comprehensive simulation studies demonstrate excellent parameter recovery, and an application to real-world data highlights the model's practical utility and enhanced predictive performance.}\\[0.4em]
{\footnotesize\bfseries Keywords:} {\footnotesize Fuzzy spatio-temporal statistics, Logistic regression, Skew random field, FMCEM algorithm. \textbackslash\{\}\textbackslash\{\}}\\[0.15em]
{\footnotesize\bfseries MSC:} {\footnotesize 62H11 , 62J12 , 62M40 . 1cm}\\
\vfill
\end{minipage}
\newpage

% p.1147
\noindent\begin{minipage}[t][\textheight]{\textwidth}
\noindent{\bfseries\large Least Angle Regression for Identifying Potential Predictors of HbA1c in Type 1 Diabetes: A Real-World Data Application}\\[0.35em]
{\itshape Elham Safari, Habibollah Esmaily, Niloofar Shabani}\\[0.2em]
{\footnotesize First author: Elham Safari}\\[0.45em]
\noindent\rule{\textwidth}{0.3pt}\\[0.4em]
{\footnotesize  Least Angle Regression (LARS) is an efficient variable selection method for high-dimensional data, offering computational advantages over traditional penalized regression approaches. This study applied LARS to predict short-term (six-month) and long-term (one-year) HbA1c levels in 271 patients with type 1 diabetes using real-world data from the Sina health information system in Mashhad, Iran. The optimal models were selected using Mallows' C\_p criterion at steps 13 and 20, respectively. Baseline HbA1c and diabetic microvascular complications were the strongest positive predictors, while dietary factors, including fruit and dairy consumption, showed protective effects. LARS provided a parsimonious, interpretable framework for identifying key potential predictors of glycemic control in type 1 diabetes using real-world clinical data.\textbackslash\{\}\textbackslash\{\}}\\[0.4em]
{\footnotesize\bfseries Keywords:} {\footnotesize Least Angle Regression(LARS), type 1 diabetes, HbA1c prediction, variable selection, and penalized regression. \textbackslash\{\}\textbackslash\{\}}\\[0.15em]
{\footnotesize\bfseries MSC:} {\footnotesize 62J07 , 62P10 , 62H30 . 1cm}\\
\vfill
\end{minipage}
\newpage

% p.1334
\noindent\begin{minipage}[t][\textheight]{\textwidth}
\noindent{\bfseries\large Order Statistics in R: Moment Computation, Goodness-of-Fit Assessment, and Statistical Inference with the mos}\\[0.35em]
{\itshape Mahdi Salehi}\\[0.2em]
{\footnotesize First author: Mahdi Salehi}\\[0.45em]
\noindent\rule{\textwidth}{0.3pt}\\[0.4em]
{\footnotesize  Although order statistics are fundamental in statistical inference, reliability analysis, and goodness-of-fit assessment, computational tools for their moment evaluation remain limited. In this paper, we briefly introduce the R package mos which provides a unified framework for the analysis of order statistics, censored samples, and k -record values. The package combines analytical evaluation of moments, whenever closed-form expressions are available, with Monte Carlo estimation for more general distributions. It also provides functions for variance, skewness, kurtosis, and efficient simulation of ordered-data structures. Applications to graphical goodness-of-fit diagnostics and best linear unbiased estimation illustrate the practical utility of the package. The mos offers a consistent and extensible environment for research and applications involving ordered data.}\\[0.4em]
{\footnotesize\bfseries Keywords:} {\footnotesize Censoring, Order statistics, Moments, Record values, Reliability, R programming language. \textbackslash\{\}\textbackslash\{\}}\\[0.15em]
{\footnotesize\bfseries MSC:} {\footnotesize 62G30 , 62 - 04 , 62N05 . 1cm}\\
\vfill
\end{minipage}
\newpage

% p.1046
\noindent\begin{minipage}[t][\textheight]{\textwidth}
\noindent{\bfseries\large Bayesian Estimation for the GIKw--Rayleigh Distribution under Various Loss Functions}\\[0.35em]
{\itshape Dr. Nahid Sanjari Fatemeh Sanatian}\\[0.2em]
{\footnotesize First author: Nahid Sanjari Fatemeh Sanatian}\\[0.45em]
\noindent\rule{\textwidth}{0.3pt}\\[0.4em]
{\footnotesize This document provides a comprehensive and self-contained treatment of Bayesian estimation for the four-parameter Generalized Inverse Kumaraswamy--Rayleigh (GIKw--RD) distribution. Starting from the probability density function, we derive the likelihood, choose independent Gamma priors, and obtain the posterior distribution. Because the posterior lacks a closed form, we present two widely used computational approaches: Lindley's approximation and Markov Chain Monte Carlo (MCMC) with Metropolis-Hastings steps. Detailed formulas are given for the Bayes estimators under Squared Error Loss (SEL), Absolute Error Loss (AEL), LINEX Loss, and General Entropy Loss (GELF). All necessary derivatives, expansions, and algorithmic steps are explicitly written out. A simulation study illustrates the performance of the methods on two generated datasets}\\[0.4em]
{\footnotesize\bfseries Keywords:} {\footnotesize Bayesian estimation, GIKw-Rayleigh distribution, Lindley approximation, MCMC, Loss functions, SEL, AEL, LINEX, GELF}\\[0.15em]
\vfill
\end{minipage}
\newpage

% p.1049
\noindent\begin{minipage}[t][\textheight]{\textwidth}
\noindent{\bfseries\large Evidential sets for the circular normal distribution}\\[0.35em]
{\itshape MohammadReza Sarvari}\\[0.2em]
{\footnotesize First author: MohammadReza Sarvari}\\[0.45em]
\noindent\rule{\textwidth}{0.3pt}\\[0.4em]
{\footnotesize  Directional data arise in various applied fields, such as wind directions, daily event times, and bird flight directions. This paper addresses evidential intervals for directional data within the evidential inference framework. To this do, a sample coming from the circular normal distribution is assumed. Intervals for the population parameters are considered when the concentration parameter is either known or unknown. For large sample sizes, some approximations are also provided to enable fast computations with massive datasets. The performance of the proposed results is illustrated by analyzing a real dataset on the times of urban injury accidents.}\\[0.4em]
{\footnotesize\bfseries Keywords:} {\footnotesize Likelihood Ratio, Circular and Directional Data, Evidential interval, Evidential sets.\textbackslash\{\}\textbackslash\{\}}\\[0.15em]
{\footnotesize\bfseries MSC:} {\footnotesize 62F25 , 62H11 , 62F15 . 1cm}\\
\vfill
\end{minipage}
\newpage

% p.1195
\noindent\begin{minipage}[t][\textheight]{\textwidth}
\noindent{\bfseries\large Semiparametric Weibull Modeling for Constant Stress Accelerated Lifetime Tests using P-spline Neural Networks}\\[0.35em]
{\itshape Bahareh Sedighi b}\\[0.2em]
{\footnotesize First author: Bahareh Sedighi b}\\[0.45em]
\noindent\rule{\textwidth}{0.3pt}\\[0.4em]
{\footnotesize  This study presents a semiparametric modeling approach for constant stress accelerated lifetime tests (CSALT) that utilizes P-spline and its artificial neural network (ANN) representation under a Type-I censoring scheme. In the proposed model, lifetimes are assumed to follow a Weibull distribution. The relationships between accelerated stresses and the characteristics of the lifetime distribution are modeled nonparametrically and additively using a B-spline approach. A penalty component is included in the loss function to regularize the model. To estimate the model parameters, a gradient-based optimization method, such as the Adam algorithm, is employed, while the smoothing parameters are estimated by minimizing the Akaike Information Criterion (AIC). The effectiveness of the proposed model is evaluated through a simulation study.}\\[0.4em]
{\footnotesize\bfseries Keywords:} {\footnotesize Artificial neural networks, Constant stress accelerated lifetime tests, P-spline, Adam algorithm\textbackslash\{\}\textbackslash\{\}}\\[0.15em]
{\footnotesize\bfseries MSC:} {\footnotesize 62N05 , 62G05 , 62N01 . 1cm}\\
\vfill
\end{minipage}
\newpage

% p.1176
\noindent\begin{minipage}[t][\textheight]{\textwidth}
\noindent{\bfseries\large A Robust Computational Framework for Time Series Clustering with Skewed Distributions}\\[0.35em]
{\itshape Ali Shadrokh}\\[0.2em]
{\footnotesize First author: Ali Shadrokh}\\[0.45em]
\noindent\rule{\textwidth}{0.3pt}\\[0.4em]
{\footnotesize  This study develops a robust model-based clustering framework for time series by incorporating Laplace and Skew-Laplace distributions to address non-normality, skewness, and heavy-tailedness in real-world data. We assume that cluster observations follow a polynomial regression structure with a temporal Vandermonde design matrix. The model parameters are estimated using both Frequentist (via the EM algorithm) and Bayesian (via MCMC sampling) approaches. Simulation results indicate that while the Skew-Laplace model performs comparably to Normal and Laplace models under symmetric conditions, it significantly outperforms them when the data are skewed and heavy-tailed. In such scenarios, symmetric distributions tend to overestimate variance to compensate for their lack of flexibility, whereas the Skew-Laplace model provides a more robust and efficient fit. Furthermore, the Bayesian approach demonstrates superior performance over Frequentist methods, particularly in clusters with small sample sizes. These findings highlight the effectiveness of the Skew-Laplace model as a comprehensive and highly adaptable framework for complex time series analysis.}\\[0.4em]
{\footnotesize\bfseries Keywords:} {\footnotesize Time series clustering, Bayesian inference, Robust estimation, Polynomial mixture models. \textbackslash\{\}\textbackslash\{\}}\\[0.15em]
{\footnotesize\bfseries MSC:} {\footnotesize 62M10 , 62H30 , 62F15 . 1cm}\\
\vfill
\end{minipage}
\newpage

% p.1100
\noindent\begin{minipage}[t][\textheight]{\textwidth}
\noindent{\bfseries\large On a Novel Weighted Multivariate Normal Distribution Using Mahalanobis Distance}\\[0.35em]
{\itshape Fatemeh Shahsanaei}\\[0.2em]
{\footnotesize First author: Fatemeh Shahsanaei}\\[0.45em]
\noindent\rule{\textwidth}{0.3pt}\\[0.4em]
{\footnotesize  The standard Multivariate Normal distribution is inherently unimodal and concentrates its probability mass at the mean, making it inadequate for modeling non-centralized or ring-shaped data. To address this limitation, we introduce the Mahalanobis-Weighted Multivariate Normal Distribution. By employing the squared Mahalanobis distance as a weighting function, the Mahalanobis-Weighted Multivariate Normal Distribution forces the central density to exactly zero, creating a topological ``hole effect'' without introducing any additional parameters. Compared to complex alternatives like Gaussian Mixture Models, the Mahalanobis-Weighted Multivariate Normal Distribution offers a strictly parsimonious and mathematically elegant solution. Simulation studies demonstrate its superior performance and computational efficiency in fitting crater-like data topologies, with broad applications in medical imaging, urban planning, and ecology.}\\[0.4em]
{\footnotesize\bfseries Keywords:} {\footnotesize Weighted Distribution, Multivariate Normal Distribution, Mahalanobis Distance, Hole Effect \textbackslash\{\}\textbackslash\{\}}\\[0.15em]
{\footnotesize\bfseries MSC:} {\footnotesize 62H05 , 60E05 , 62E15 . 1cm}\\
\vfill
\end{minipage}
\newpage

% p.1009
\noindent\begin{minipage}[t][\textheight]{\textwidth}
\noindent{\bfseries\large The Ising spin model on the square grid system based on Boltzmann distribution}\\[0.35em]
{\itshape Mehdi Shams, Mohammad Ali Mirzaie}\\[0.2em]
{\footnotesize First author: Mehdi Shams}\\[0.45em]
\noindent\rule{\textwidth}{0.3pt}\\[0.4em]
{\footnotesize  Many scientific disciplines have adopted the Ising model as a fundamental physical framework. Three specific configurations of this model are examined in this paper, with emphasis on the square n n grid. For general Ising models built upon the Boltzmann stationary distribution, we show that the maximum likelihood estimator coincides with the method-of-moments estimator for the inverse temperature parameter . Analytical solution of the resulting equation proves impossible; hence linear regression methods are employed to obtain these estimates. A general condition governing convexity of the Boltzmann probability entropy is also put forward.}\\[0.4em]
{\footnotesize\bfseries Keywords:} {\footnotesize Entropy, Boltzmann distribution, maximum likelihood estimator, regression models, induced probability vector. \textbackslash\{\}\textbackslash\{\}}\\[0.15em]
{\footnotesize\bfseries MSC:} {\footnotesize 62B10 , 62F10 , 62P35 , 62M10 . 1cm}\\
\vfill
\end{minipage}
\newpage

% p.1183
\noindent\begin{minipage}[t][\textheight]{\textwidth}
\noindent{\bfseries\large Time-Dependent Cumulative Past Fisher Information: Bounds, Orderings, and Mixture Representations}\\[0.35em]
{\itshape Maryam Sharafi}\\[0.2em]
{\footnotesize First author: Maryam Sharafi}\\[0.45em]
\noindent\rule{\textwidth}{0.3pt}\\[0.4em]
{\footnotesize  This paper develops dynamic versions of cumulative past Fisher information measures for the past lifetime and inactivity time of a system. The proposed information functionals evolve over time and capture uncertainty from a retrospective viewpoint. We express them in terms of the reversed hazard rate and the mean inactivity time, derive sharp lower bounds using Cauchy--Schwarz and convexity arguments, and link them to reliability aging notions such as decreasing mean inactivity time and log-convexity. Stochastic orders based on the new dynamic information measure are introduced and characterized. Explicit formulas are obtained for the proportional reversed hazards model and for arithmetic mixture models, where the Fisher information is shown to equal four times a power mean of chi-square distances. Nonparametric estimators are also provided.\textbackslash\{\}\textbackslash\{\}}\\[0.4em]
{\footnotesize\bfseries Keywords:} {\footnotesize Cumulative past entropy, Mean inactivity time, Mixture models, Reversed hazard rate, Stochastic orders. \textbackslash\{\}\textbackslash\{\}}\\[0.15em]
{\footnotesize\bfseries MSC:} {\footnotesize 62N05 , 94A17 , 62B10 . 1cm}\\
\vfill
\end{minipage}
\newpage

% p.1198
\noindent\begin{minipage}[t][\textheight]{\textwidth}
\noindent{\bfseries\large Asymptotic Power Analysis for Testing Points of Impact under Brownian Motion Predictors}\\[0.35em]
{\itshape Alireza Shirvani}\\[0.2em]
{\footnotesize First author: Alireza Shirvani}\\[0.45em]
\noindent\rule{\textwidth}{0.3pt}\\[0.4em]
{\footnotesize  Functional linear regression models with points of impact provide an important framework for identifying localized effects of functional predictors. In this paper, the asymptotic behavior of the test statistic for points of impact, proposed by Shirvani2024 , is investigated under Brownian motion predictors. Using the explicit spectral decomposition of Brownian motion, simplified expressions for the covariance structure and the asymptotic power function are derived. Under local alternatives, the limiting distribution of the test statistic is shown to converge to a noncentral chi-square distribution with an explicit noncentrality parameter. Simulation studies demonstrate satisfactory finite-sample performance and illustrate the effect of the impact location on the asymptotic power of the proposed test.}\\[0.4em]
{\footnotesize\bfseries Keywords:} {\footnotesize Functional data analysis, Functional linear regression, Points of impact, Brownian motion, Asymptotic power. \textbackslash\{\}\textbackslash\{\}}\\[0.15em]
{\footnotesize\bfseries MSC:} {\footnotesize 62G35 , 62J05 , 60G15 . 1cm}\\
\vfill
\end{minipage}
\newpage

% p.1217
\noindent\begin{minipage}[t][\textheight]{\textwidth}
\noindent{\bfseries\large Jackknife Ridge Estimation in Linear Mixed Models with Measurement Error}\\[0.35em]
{\itshape Mina Shirvani, Abdolrahman Rasekh, Babak Babadi}\\[0.2em]
{\footnotesize First author: Mina Shirvani}\\[0.45em]
\noindent\rule{\textwidth}{0.3pt}\\[0.4em]
{\footnotesize  This study investigates the problem of bias in linear mixed models (LMMs) when ridge estimation is applied in the presence of measurement errors in fixed-effect variables. First, the ridge estimator for LMMs is described, and then a new estimator, called the jackknife ridge estimator (JRE), is proposed. The performance of the proposed estimator is compared with the conventional ridge estimator, demonstrating improvements in terms of bias and mean squared error. In addition, the asymptotic properties of both estimators are derived. Finally, simulation studies and a numerical example are presented to evaluate the efficiency of the proposed estimator.}\\[0.4em]
{\footnotesize\bfseries Keywords:} {\footnotesize Jackknife ridge; Linear mixed models; Measurement error; Bias reduction; Asymptotic properties \textbackslash\{\}\textbackslash\{\}}\\[0.15em]
{\footnotesize\bfseries MSC:} {\footnotesize 99X99 , 99X99 , 99X99 . 1cm}\\
\vfill
\end{minipage}
\newpage

% p.1057
\noindent\begin{minipage}[t][\textheight]{\textwidth}
\noindent{\bfseries\large A Novel Approach for Statistical Shape Analysis Using the Square Root Velocity Function Representation}\\[0.35em]
{\itshape A Novel Approach for Statistical Shape Analysis Using SRVF}\\[0.2em]
{\footnotesize First author: A Novel Approach for Statistical Shape Analysis Using SRVF}\\[0.45em]
\noindent\rule{\textwidth}{0.3pt}\\[0.4em]
{\footnotesize  abstract Statistical shape analysis for manifold-valued data poses significant challenges in comparing shapes and modeling their variability. Classical methods based on tangent space principal component analysis (TPCA) rely on local linearization and often suffer from limitations in stability and interpretability. In this work, we introduce a novel framework for functional principal component analysis (FPCA) in the Sobolev space H\textasciicircum{}1([0,1]) , which naturally incorporates geometric information of curves. Shapes are represented using the square root velocity function (SRVF) and aligned via the Karcher mean to remove phase variability. The variability is then analyzed globally in H\textasciicircum{}1 , where the covariance operator accounts for both the functions and their derivatives. We clarify the distinct advantages of this global approach over local TPCA approximations and standard L\textasciicircum{}2 -based FPCA. We show theoretically that the resulting principal components are equivalent, up to first-order variability, to those obtained from tangent-space approaches. Experiments on real two-dimensional contour datasets demonstrate that the proposed method captures most of the structural variability with few components while maintaining consistent numerical stability and reconstruction accuracy. abstract}\\[0.4em]
{\footnotesize\bfseries Keywords:} {\footnotesize Statistical Shape Analysis, Square Root Velocity Function (SRVF), Functional Principal Component Analysis (FPCA), Banach Space H\textasciicircum{}1 , Covariance Operator.\textbackslash\{\}\textbackslash\{\}}\\[0.15em]
{\footnotesize\bfseries MSC:} {\footnotesize 62R30 , 62H25 . \textbackslash\{\}\textbackslash\{\}}\\
\vfill
\end{minipage}
\newpage

% p.1300
\noindent\begin{minipage}[t][\textheight]{\textwidth}
\noindent{\bfseries\large Mixed-Effects Hurdle Trees for Longitudinal Zero-Inflated Count Data}\\[0.35em]
{\itshape Elham Tabrizi}\\[0.2em]
{\footnotesize First author: Elham Tabrizi}\\[0.45em]
\noindent\rule{\textwidth}{0.3pt}\\[0.4em]
{\footnotesize  Longitudinal count data with excess zeros are common in many applied fields, yet flexible tree-based methods for such data remain scarce. This paper introduces the Penalized Quasi-Likelihood Hurdle Tree (PHT), a novel algorithm that combines regression trees with mixed-effects models to capture nonlinear fixed effects and subject specific correlations in zero-inflated longitudinal counts. The PHT employs a penalized quasi-likelihood approximation, alternating between estimating random effects and updating tree-based fixed effects via weighted pseudo-responses. Simulation studies demonstrate that PHT outperforms standard regression trees, Poisson GLMMs, and hurdle GLMMs in terms of predictive accuracy. An application to a salamander abundance dataset further confirms the practical utility of PHT, revealing key ecological drivers while achieving the lowest out-of-sample prediction error.}\\[0.4em]
{\footnotesize\bfseries Keywords:} {\footnotesize Longitudinal data, Zero-inflated count data, Hurdle model, Mixed-effects model, Regression tree, Penalized quasi-likelihood.\textbackslash\{\}\textbackslash\{\}}\\[0.15em]
{\footnotesize\bfseries MSC:} {\footnotesize 99X99 , 99X99 , 99X99 . 1cm}\\
\vfill
\end{minipage}
\newpage

% p.1189
\noindent\begin{minipage}[t][\textheight]{\textwidth}
\noindent{\bfseries\large Parameters Estimation of the Lomax-Lindley Distribution Utilizing Record Data}\\[0.35em]
{\itshape Bahman Tarvirdizade}\\[0.2em]
{\footnotesize First author: Bahman Tarvirdizade}\\[0.45em]
\noindent\rule{\textwidth}{0.3pt}\\[0.4em]
{\footnotesize  This paper is devoted to the problem of parameters estimation based on record values from Lomax-Lindley distribution. The estimation of parameters are obtained by using the maximum likelihood method and Bayesian method under squared error and linear-exponential loss functions. A Monte Carlo simulation study is conducted to investigate and compare the performance of different types of estimators presented in this paper. Based on our simulation outcomes, it is observed that the Bayesian estimators perform better than the MLEs in the estimation of the unknown parameters.}\\[0.4em]
{\footnotesize\bfseries Keywords:} {\footnotesize Lomax-Lindley distribution, Maximum likelihood estimation, Bayesian estimation, and Record values. \textbackslash\{\}\textbackslash\{\}}\\[0.15em]
{\footnotesize\bfseries MSC:} {\footnotesize 62F15 , 65C05 . 1cm}\\
\vfill
\end{minipage}
\newpage

% p.1034
\noindent\begin{minipage}[t][\textheight]{\textwidth}
\noindent{\bfseries\large Bivariate Quantile Regression Analysis for Longitudinal Cardiovascular Risk Factors}\\[0.35em]
{\itshape Fahimeh Tourani-Farani}\\[0.2em]
{\footnotesize First author: Fahimeh Tourani-Farani}\\[0.45em]
\noindent\rule{\textwidth}{0.3pt}\\[0.4em]
{\footnotesize  Modeling the joint evolution of correlated cardiovascular risk factors across their full distribution remains a statistical challenge, as traditional mean-based longitudinal methods fail to capture extreme-dependent associations and heterogeneous covariate effects. We present a bivariate quantile regression framework for longitudinal data based on the generalized Johnson distribution and the Rosenblatt transformation. The method enables simultaneous modeling of two correlated outcomes at arbitrary quantile levels while preserving interpretability and automatically avoiding quantile crossing. A computationally efficient maximum likelihood estimator is developed with consistency and asymptotic normality guarantees. A comprehensive simulation study demonstrates that the proposed quantile regression model consistently outperforms all competing methods across all evaluation metrics.}\\[0.4em]
{\footnotesize\bfseries Keywords:} {\footnotesize Bivariate analysis, Correlated outcomes, Generalized Johnson distribution, Longitudinal data, Rosenblatt transformation. \textbackslash\{\}\textbackslash\{\}}\\[0.15em]
{\footnotesize\bfseries MSC:} {\footnotesize 99X99 , 99X99 , 99X99 . 1cm}\\
\vfill
\end{minipage}
\newpage

% p.1041
\noindent\begin{minipage}[t][\textheight]{\textwidth}
\noindent{\bfseries\large European option valuation with modified Heston model under stochastic market liquidity}\\[0.35em]
{\itshape Parsa Yahyavi, p}\\[0.2em]
{\footnotesize First author: Parsa Yahyavi}\\[0.45em]
\noindent\rule{\textwidth}{0.3pt}\\[0.4em]
{\footnotesize  Comprehensive modeling of stock prices plays a critical role in financial derivatives valuation. Although, the classic Heston model captures the stochastic volatility, it struggles to replicate the term structure of implied volatility and assumes a perfectly liquid market. To address these inconsistencies, we introduce a Heston model by incorporating independent Ornstein--Uhlenbeck processes for both the long-term mean of volatility and stochastic market liquidity. Moreover, we derive an analytical European option pricing formula and validate the model using real market data from leading companies. We employ the root mean square error to assess the superiority of the introduced model against three benchmark models in replicating European option prices and implied volatility.}\\[0.4em]
{\footnotesize\bfseries Keywords:} {\footnotesize European option pricing, Heston model, Stochastic liquidity, Stochastic volatility \textbackslash\{\}\textbackslash\{\}}\\[0.15em]
{\footnotesize\bfseries MSC:} {\footnotesize 91G20 , 91G30 , 60J76 . 1cm}\\
\vfill
\end{minipage}
\newpage

% p.1084
\noindent\begin{minipage}[t][\textheight]{\textwidth}
\noindent{\bfseries\large Stress--Strength Reliability with Clayton Copula and Generalized Rayleigh Marginals}\\[0.35em]
{\itshape Yarahmadi et al}\\[0.2em]
{\footnotesize First author: Yarahmadi}\\[0.45em]
\noindent\rule{\textwidth}{0.3pt}\\[0.4em]
{\footnotesize  This study estimates stress-strength reliability when the stress and strength variables are dependent. The dependence is modeled by the Clayton copula, which captures lower-tail association, and the marginals follow the generalized Rayleigh distribution. A unified framework is developed that includes maximum likelihood estimation, bootstrap confidence intervals, and Bayesian estimation under different loss functions using a Metropolis-within-Gibbs sampler. A Monte Carlo simulation evaluates the performance of the estimators, and a real data example shows the practical use of the model. The findings indicate that ignoring dependence can lead to biased reliability estimates.}\\[0.4em]
{\footnotesize\bfseries Keywords:} {\footnotesize Generalized Rayleigh distribution, Clayton copula, stress-strength reliability, Bayesian inference. \textbackslash\{\}\textbackslash\{\}}\\[0.15em]
{\footnotesize\bfseries MSC:} {\footnotesize 62N05 , 62F15 , 62H05 . 1cm}\\
\vfill
\end{minipage}
\newpage

% p.1162
\noindent\begin{minipage}[t][\textheight]{\textwidth}
\noindent{\bfseries\large Comparative Evaluation of Synthetic Medical Data Generation Methods in Pediatric Osteopenia}\\[0.35em]
{\itshape Razieh Yousefi}\\[0.2em]
{\footnotesize First author: Razieh Yousefi}\\[0.45em]
\noindent\rule{\textwidth}{0.3pt}\\[0.4em]
{\footnotesize  Medical datasets, often suffer from limited sample sizes. Synthetic data generation may support data augmentation, privacy-preserving sharing, and methodological validation. This study compared six synthetic data generation strategies using a pediatric osteopenia dataset ( n=71 ) from children with type 1 diabetes mellitus, including conditional and non-conditional Bootstrap, Multivariate Normal (MVN), and Gaussian Copula methods. Synthetic datasets of 500 samples were evaluated using utility (TSTR ROC analysis), fidelity (Kolmogorov--Smirnov testing and correlation preservation), and privacy (Nearest Neighbor Distance Ratio; NNDR) assessments. The real-data model achieved an AUC of 0.796 . Conditional Bootstrap ( AUC = 0.786 ), conditional MVN ( AUC = 0.776 ), and conditional Copula ( AUC = 0.776 ) preserved predictive performance well, whereas the non-conditional Copula approach performed poorly ( AUC = 0.224 , p = 0.011 ). Bootstrap and Copula methods showed strong distributional similarity, while conditional MVN demonstrated the best correlation preservation. Privacy analysis indicated high disclosure risk for Bootstrap methods, whereas MVN and Copula approaches provided substantially better privacy protection. Overall, conditional MVN provided the best balance between utility, fidelity, and privacy preservation in this low-sample clinical setting.}\\[0.4em]
{\footnotesize\bfseries Keywords:} {\footnotesize Synthetic data, Multivariate normal, Gaussian copula, Bootstrap, Pediatric osteopenia}\\[0.15em]
{\footnotesize\bfseries MSC:} {\footnotesize 62H12 , 62P10 , 68T07 . 1cm}\\
\vfill
\end{minipage}
\newpage

% p.1052
\noindent\begin{minipage}[t][\textheight]{\textwidth}
\noindent{\bfseries\large Likelihood-Based Evidential Comparison of Classical Goodness-of-Fit Tests}\\[0.35em]
{\itshape Fatemeh Yousefzadeh}\\[0.2em]
{\footnotesize First author: Fatemeh Yousefzadeh}\\[0.45em]
\noindent\rule{\textwidth}{0.3pt}\\[0.4em]
{\footnotesize  The Law of Likelihood establishes that the strength of statistical evidence favoring one hypothesis over another is represented by their likelihood ratio. Within the framework of evidential statistics, this study compares the performance of several classical goodness-of-fit tests---including the Anderson--Darling, Kolmogorov--Smirnov, Cramér--von Mises, Watson, and Likelihood Ratio Test---based on their ability to yield reliable and interpretable evidence. Simulation results across different sample sizes quantify the probabilities of weak and misleading evidence under both null and alternative models, offering new insights into the evidential performance of these tests.}\\[0.4em]
{\footnotesize\bfseries Keywords:} {\footnotesize Evidential statistics, Goodness-of-fit tests, Law of Likelihood, Misleading evidence, Weak evidence. \textbackslash\{\}\textbackslash\{\}}\\[0.15em]
{\footnotesize\bfseries MSC:} {\footnotesize 62F25 , 62G15 . 1cm}\\
\vfill
\end{minipage}
\newpage

% p.1180
\noindent\begin{minipage}[t][\textheight]{\textwidth}
\noindent{\bfseries\large EndSurvGAN: End-to-End Survival-Aware Generative Modeling of Censored Time-to-Event Data}\\[0.35em]
{\itshape Z.Esfandiar}\\[0.2em]
{\footnotesize First author: Z.Esfandiar}\\[0.45em]
\noindent\rule{\textwidth}{0.3pt}\\[0.4em]
{\footnotesize Generating realistic synthetic survival data remains a challenging task due to the complex dependencies among covariates, survival times, and censoring indicators. Existing approaches often rely on stage-wise frameworks that separately model feature generation and survival outcomes, potentially weakening the coherence of the underlying joint survival distribution. In this paper, we propose EndSurvGAN , an end-to-end survival-aware generative adversarial network designed to jointly model covariates, survival times, and censoring indicators within a unified framework. By directly learning the joint survival distribution, the proposed model preserves both survival dynamics and censoring mechanisms during data generation. The proposed framework is evaluated on three benchmark survival datasets, namely SUPPORT, ACTG, and Rotterdam. Experimental results indicate improved distributional fidelity, better preservation of Kaplan Meier survival curves, and enhanced downstream predictive utility compared with existing tabular and survival-specific generative models. These findings suggest that end-to-end survival-aware generative modeling provides an effective framework for synthetic censored survival data generation.}\\[0.4em]
{\footnotesize\bfseries Keywords:} {\footnotesize Survival Analysis, Synthetic Data Generation, Generative Adversarial Networks, Censored Survival Data, Deep Survival Learning, Healthcare Analytics \textbackslash\{\}\textbackslash\{\}}\\[0.15em]
{\footnotesize\bfseries MSC:} {\footnotesize 99X99 , 99X99 , 99X99 . 1cm}\\
\vfill
\end{minipage}
\newpage

% p.1005
\noindent\begin{minipage}[t][\textheight]{\textwidth}
\noindent{\bfseries\large Stein Estimators in Right-Censored Zero-inflated Bell Models}\\[0.35em]
{\itshape Zahra Zandi}\\[0.2em]
{\footnotesize First author: Zahra Zandi}\\[0.45em]
\noindent\rule{\textwidth}{0.3pt}\\[0.4em]
{\footnotesize  In the current paper, we propose Stein and positive Stein estimators to improve the efficiency of parameters estimation in the zero-inflated Bell models with right-censored data, assuming that some predictor variables are non-significant. A Monte Carlo simulation study is conducted under various settings, including different censoring constants, and non-significant predictors, to compare the performance of our estimators with the unrestricted estimator. The results demonstrate the superiority of our approach.}\\[0.4em]
{\footnotesize\bfseries Keywords:} {\footnotesize Right-censored data, Stein estimator, Positive Stein estimator, Zero-inflated Bell model. \textbackslash\{\}\textbackslash\{\}}\\[0.15em]
{\footnotesize\bfseries MSC:} {\footnotesize 62J07 , 62N01 . 1cm}\\
\vfill
\end{minipage}
\newpage

% p.1066
\noindent\begin{minipage}[t][\textheight]{\textwidth}
\noindent{\bfseries\large p.1066}\\[0.35em]
{\itshape }\\[0.2em]
{\footnotesize First author: }\\[0.45em]
\noindent\rule{\textwidth}{0.3pt}\\[0.4em]
{\footnotesize }\\[0.4em]
\vfill
\end{minipage}
\newpage

% p.1333
\noindent\begin{minipage}[t][\textheight]{\textwidth}
\noindent{\bfseries\large An unified view of probabilistic inference under incomplete information: Maximum Entropy and Bayesian Updating}\\[0.35em]
{\itshape Mathematics Subject Classification (2010 :}\\[0.2em]
{\footnotesize First author: }\\[0.45em]
\noindent\rule{\textwidth}{0.3pt}\\[0.4em]
{\footnotesize  Statistical inference under incomplete information involves two complementary but distinct tasks: assigning an initial probability distribution when knowledge is limited to partial constraints, and updating that distribution when observations become available. Maximum Entropy addresses the first task by selecting a distribution consistent with the known moment conditions, thereby avoiding assumptions unsupported by the available information. Bayesian inference addresses the second by using the likelihood to revise prior probabilities in light of observed data. Recent work connects these roles within a single Bayesian framework: moment constraints determine the prior weighting of candidate distributions, while observations modify those weights through the likelihood. When the constraint-based prior is sharply concentrated around the entropy-maximizing distribution, the Bayesian result approaches the Maximum Entropy solution. This review examines this connection and illustrates the resulting probabilistic updating process through a numerical experiment.}\\[0.4em]
{\footnotesize\bfseries Keywords:} {\footnotesize Maximum Entropy, Bayesian Inference, and Entropic Priors. \textbackslash\{\}\textbackslash\{\}}\\[0.15em]
{\footnotesize\bfseries MSC:} {\footnotesize 99X99 , 1cm}\\
\vfill
\end{minipage}
\newpage

% p.1096
\noindent\begin{minipage}[t][\textheight]{\textwidth}
\noindent{\bfseries\large Application of Best Linear Unbiased Prediction (BLUP in Stability Analysis and Genotype × Environment Interaction Studies}\\[0.35em]
{\itshape }\\[0.2em]
{\footnotesize First author: }\\[0.45em]
\noindent\rule{\textwidth}{0.3pt}\\[0.4em]
{\footnotesize Genotype × environment interaction (GEI) is a major challenge in plant breeding, as it complicates the identification of genotypes that are both high-performing and stable across diverse environments. Multi-environment trials (METs) generate complex datasets, often unbalanced and incomplete, where traditional ANOVA or regression-based stability metrics are limited. Best Linear Unbiased Prediction (BLUP), within a mixed model framework, provides a robust approach for estimating genotypic performance while accounting for random effects, environmental variation, and heterogeneous residuals. BLUP enables shrinkage-based predictions, improving accuracy over raw means or BLUE, particularly in unbalanced datasets. Furthermore, incorporation of pedigree or marker-derived relationship matrices allows exploitation of genetic correlations, enhancing prediction of untested genotypes and facilitating genomic selection. BLUP-based methods can also quantify stability through variance of predicted genotypic values across environments, factor-analytic decomposition of G×E, and reliability-adjusted indices, providing actionable insights for selection. This manuscript reviews theoretical foundations, practical implementation, and computational approaches for BLUP in METs, including hierarchical and factorial designs. It also discusses integration with factor-analytic models for complex GEI, the use of combined pedigree and marker information, and strategies for handling missing data. BLUP-based stability analysis represents a flexible and powerful tool for modern plant breeding, enabling selection of widely adapted and specifically adapted genotypes while maximizing genetic gain under diverse environmental conditions.}\\[0.4em]
{\footnotesize\bfseries Keywords:} {\footnotesize Mixed models, Plant breeding, Multi-environment trials}\\[0.15em]
{\footnotesize\bfseries MSC:} {\footnotesize 62P10, 62J12, 62H25, 92D10. introduction}\\
\vfill
\end{minipage}
\newpage

% p.1097
\noindent\begin{minipage}[t][\textheight]{\textwidth}
\noindent{\bfseries\large Effective Analysis of Genetic Diversity of Plant Species through Artificial Intelligence Tools}\\[0.35em]
{\itshape }\\[0.2em]
{\footnotesize First author: }\\[0.45em]
\noindent\rule{\textwidth}{0.3pt}\\[0.4em]
{\footnotesize Genetic variation within species arises from complex interactions between evolutionary processes and environmental factors, often reflected in polymorphic patterns of Single-Nucleotide Polymorphisms (SNPs). Traditional landscape genetics approaches have primarily focused on either broad-scale or local-scale environmental effects, limiting the statistical inference of genetic diversity. In this study, we present a comprehensive statistical framework integrating supervised machine learning, convolutional neural networks (CNNs), and simulation-based modeling to quantify the influence of multiple factors on genetic variation. Supervised machine learning allows prediction of key genetic metrics, such as allelic richness and population differentiation, while simultaneously estimating the relative importance of environmental and historical drivers. CNNs process SNP alignments as high-dimensional data, extracting patterns without requiring predefined summary statistics or population assignments. Simulation-based approaches, including Approximate Bayesian Computation (ABC), enable robust estimation of parameter posteriors under complex evolutionary scenarios. By combining these methods, our framework captures both discrete and continuous genetic processes, models stochasticity in natural populations, and provides interpretable statistical measures of factor importance. This approach enhances inference across spatial and temporal scales, facilitating hypothesis testing and evolutionary interpretation. The framework also has practical applications for conservation planning and biodiversity management, allowing researchers to identify populations and environmental factors most critical for maintaining genetic diversity. This integrative, statistics-driven approach represents a significant advancement in molecular ecology, offering scalable, interpretable, and robust tools for analyzing genomic variation in natural populations.}\\[0.4em]
{\footnotesize\bfseries Keywords:} {\footnotesize Supervised machine learning, Convolutional neural networks, Statistical genomics}\\[0.15em]
{\footnotesize\bfseries MSC:} {\footnotesize 62P10, 62F15, 92D20, 92D15. introduction}\\
\vfill
\end{minipage}
\newpage

% p.1315
\noindent\begin{minipage}[t][\textheight]{\textwidth}
\noindent{\bfseries\large Feature Ranking and Incremental Selection for High-Dimensional Classification}\\[0.35em]
{\itshape Mathematics Subject Classification (2010 :}\\[0.2em]
{\footnotesize First author: }\\[0.45em]
\noindent\rule{\textwidth}{0.3pt}\\[0.4em]
{\footnotesize  High-dimensional datasets in genomics pose classification challenges due to the curse of dimensionality and redundant features. Feature selection addresses this by identifying informative subsets to improve model interpretability, efficiency, and performance. However, many existing feature selection methods lack stability, meaning that small changes in the data can lead to drastically different selected feature subsets. This paper introduces an integrated stability-ranked feature selection (ISRS) framework designed to address this limitation. ISRS combines ensemble-based feature ranking with a novel stability metric to identify a predictive feature subset. We evaluate ISRS on two real datasets, demonstrating its superior stability and competitive or improved classification accuracy compared to established regularized methods like Lasso and Elastic Net.}\\[0.4em]
{\footnotesize\bfseries Keywords:} {\footnotesize Feature ranking, Feature selection, Selection stability, Support vector machine. \textbackslash\{\}\textbackslash\{\}}\\[0.15em]
{\footnotesize\bfseries MSC:} {\footnotesize 62J07 , 62F07 , 62H30 . 1cm}\\
\vfill
\end{minipage}
\newpage

% p.1073
\noindent\begin{minipage}[t][\textheight]{\textwidth}
\noindent{\bfseries\large Maximum likelihood estimation of length-biased censored data under Frank copula dependency}\\[0.35em]
{\itshape Mathematics Subject Classification (2010 :}\\[0.2em]
{\footnotesize First author: }\\[0.45em]
\noindent\rule{\textwidth}{0.3pt}\\[0.4em]
{\footnotesize  Ignoring censoring causes bias in survival analysis, especially for length-biased data. In this work, we estimate parameters in length-biased survival data with dependent censoring using a copula function to model the dependency between survival and censoring times. Extensive simulations show that our approach improves estimation accuracy and simplifies computing conditional probabilities of copulas. A real-data application is also provided.\textbackslash\{\}\textbackslash\{\}}\\[0.4em]
{\footnotesize\bfseries Keywords:} {\footnotesize Survival analysis, Length-biased, Copula, Parameter estimation \textbackslash\{\}\textbackslash\{\}}\\[0.15em]
{\footnotesize\bfseries MSC:} {\footnotesize 62H05 , 62Hxx , 62N02 . 1cm}\\
\vfill
\end{minipage}
\newpage

% p.1094
\noindent\begin{minipage}[t][\textheight]{\textwidth}
\noindent{\bfseries\large Method for Estimating Genotype Stability, Adaptability, and Environmental Discrimination in Plant Breeding Programs}\\[0.35em]
{\itshape }\\[0.2em]
{\footnotesize First author: }\\[0.45em]
\noindent\rule{\textwidth}{0.3pt}\\[0.4em]
{\footnotesize Breeding programs aim not only to improve crop performance but also to ensure stability and adaptability across diverse environments. Evaluating genotype resilience requires understanding both general adaptability ability (GAA), the mean trait performance across multiple environments, as well as specific adaptability ability (SAA), the deviation of a genotype in a particular environment. A comprehensive approach for assessing GAA, SAA, and genotype stability, is presented for integrating both linear and nonlinear responses to environmental variation. The method relies on multi-environment trials, partitioning phenotypic variance into genotype, environment, and genotype × environment interaction components using ANOVA models, and calculating stability indices such as the relative stability and stability value of genotypes (SVG). This framework allows breeders to select genotypes with optimal performance while maintaining stability, balancing trade-offs between high mean trait values and environmental sensitivity. Furthermore, it provides insights into environmental differentiation and the compensatory or destabilizing effects of genotype × environment interactions. By combining theoretical and practical aspects of adaptive breeding, this approach offers a robust tool for guiding selection strategies that enhance both productivity and resilience in crop populations. To compare populations and select appropriate breeding methods across different environments, a technique is introduced to partition the phenotypic variance of a population into GAA and SAA variances.}\\[0.4em]
{\footnotesize\bfseries Keywords:} {\footnotesize Adaptability ability, Stability value, Genotype × environment interaction}\\[0.15em]
{\footnotesize\bfseries MSC:} {\footnotesize 62P10, 62J10, 62K10, 92D10. introduction}\\
\vfill
\end{minipage}
\newpage

\end{document}